% Información para el motor (arara). Compilar en LuaLaTeX o XeLaTeX sólamente.

%% arara: lualatex: {
%% arara: --> interaction: nonstopmode,
%% arara: --> options: ['-file-line-error'],
%% arara: --> shell: yes,
%% arara: --> synctex: yes,
%% arara: --> }

%arara: lualatex: {
%arara: --> interaction: nonstopmode,
%arara: --> options: ['-file-line-error'],
%arara: --> shell: yes,
%arara: --> synctex: yes,
%arara: --> }
%arara: bib2gls


%% arara: makeindex

\documentclass[letterpaper,twoside]{article}

 %%%%%%% VARIABLES ESTÁTICAS %%%%%%%
% Aqui se definen las variables para diferentes procedimientos.
% Verificar que las unidades coincidan si se deciden cambiar.

\newcommand{\defglo}[1]{\item[{\bfseries \glsname{#1}}] \glsdesc{#1}} % Define un comando para dar un item en un entorno de descripción para una palabra definida en el glosario
% \includeonly{src/PSA-2.19.tex,src/PSA-Gloss.tex}
\newcommand{\TipoID}{TEMP} % Definirlo como: PRO, PROG, ToC, G, L, CE, CI, I, etc.
\newcommand{\Prog}{TEMP}
\newcommand{\ProgTitulo}{TEMP}

\newcommand{\ie}{\emph{i.e.}}
\newcommand{\eg}{\emph{e.g.}}

%%%%%%%%%%%%%%%% Formularios %%%%%%%%%%%%%%%%
\newcommand{\IdFormAACC}{F-OP-40}

\newcommand{\BitES}{Bitácora de entradas y salidas --- Clientes (F-AD-1)}
\newcommand{\BitULl}{Bitácora de uso de llaves (F-AC-1)}
\newcommand{\BitVig}{Bitácora de vigilancia (F-AC-2)}
\newcommand{\Fpap}{Papeleta (F-AC-3)}
\newcommand{\Oent}{Orden de entrada (F-AC-4)}
\newcommand{\Osal}{Orden de salida (F-AC-5)}
\newcommand{\LayES}{Layout de carga o descarga (F-AC-6)}
\newcommand{\BitClor}{Bitácora de cloración de agua (F-AC-7)}
\newcommand{\BitTap}{Bitácora de concentración de tapete sanitario (F-AC-8)}
\newcommand{\BitUQ}{Bitácora de uso de productos químicos (F-AC-9)}
\newcommand{\BitLAd}{Bitácora de limpieza de aduana sanitaria (F-AC-10)}
\newcommand{\BitLOf}{Bitácora de limpieza de oficinas (F-AC-11)}
\newcommand{\BitLEx}{Bitácora de limpieza de exteriores (F-AC-12)}
\newcommand{\BitBPH}{Bitácora de cumplimiento de BPH (F-HyS-1)}
\newcommand{\BitCAP}{Bitácora de consumo de agua potable (F-HyS-2)}
\newcommand{\BitCE}{Bitácora de consumo de energía eléctrica (F-HyS-3)}
\newcommand{\BitIEx}{Bitácora de inspección de extintores (F-HyS-4)}
\newcommand{\BitVIR}{Bitácora de verificación de termómetros infrarrojos (F-HyS-5)}
\newcommand{\LVCMo}{Lista de verificación de condición de montacargas (F-HyS-6)}
\newcommand{\RTAC}{Registro de temperaturas con termómetro IR --- Aseguramiento de calidad (F-HyS-7)}
\newcommand{\RTMa}{Registro de temperaturas con termómetro IR --- Mantenimiento (F-HyS-8)}
\newcommand{\RAcL}{Registro de actividad de limpieza (F-HyS-9)}
\newcommand{\RAcM}{Registro de actividad de mantenimiento (F-Hys-10)}
\newcommand{\ReEq}{Registro de revisión de equipos de enfriamiento (F-MAN-1)}
\newcommand{\RCMi}{Registro de inspección de infraestructura en prevención (F-MAN-2)}
\newcommand{\Rcap}{Registro de asistencia a capacitación (F-MAN-3)}
\newcommand{\REvDQ}{Registro de eventos de derrame de químicos (F-MAN-4)}
\newcommand{\InvVyPD}{Registro de inventario de vidrio y plástico duro (F-MAN-5)}
\newcommand{\RVerSGC}{Registro de verificación del sistema de gestión de calidad (F-MAN-6)}
\newcommand{\RRCA}{Registro de recorrido de control de alérgenos (F-MC-1)}
\newcommand{\RAC}{Registro de acción correctiva o preventiva (F-OP-1)}
\newcommand{\RACQ}{Registro de acción correctiva --- Quejas (F-OP-2)}
\newcommand{\RVSCP}{Registro de verificación del servicio de control de plagas (F-OP-3)}
\newcommand{\RevD}{Registro de revisión diaria (F-SEG-1)}
\newcommand{\RCVyPD}{Registro de evento de contaminación por VyPD (F-SEG-2)}
\newcommand{\REvCA}{Registro de evento de contaminación por alérgeno (F-SEG-3)}



%%%%%%%%%%%%%%%% CÓDIGOS PSA %%%%%%%%%%%%%%%%
% \newcommand{\PSA}{PSA-3-PRO}                 % PSA-2.3 Inspección de producto durante su almacenamiento

%%%%%%%%%%%%%%%% PSA-2.3 | Inspección de producto durante su almacenamiento %%%%%%%%%%%%%%%%
\newcommand{\TiempoAndenCong}{\qty{\leq 30}{\minute}}
\newcommand{\TiempoAndenRefri}{\qty{\leq 30}{\minute}}
\newcommand{\TiempoAveriaRefri}{4}              % Tiempo máximo de permanencia de alimento refrigerado en cámara averiada (en horas) 
\newcommand{\TiempoAveriaConge}{10}             % Tiempo máximo de permanencia de alimento congelado en cámara averiada (en horas)
\newcommand{\EspacioMinimoParedProducto}{30}    % Espacio entre tarima y pared (en centimetros)
\newcommand{\EspacioMinimoTechoProducto}{30}    % Espacio entre tarima y techo (en centimetros)
\newcommand{\VecesTempManual}{dos}              % Veces que se verifica la temperatura manualmente



%%%%%%%%%%%%%%%% PSA-2.3 | Política de aprobación de proveedores %%%%%%%%%%%%%%%%
\newcommand{\NotifAuditoriaAProveedor}{2 semanas}           % Plazo máximo para notificar a proveedor sobre auditoría.
\newcommand{\PlazoResultadoAProveedor}{2 semanas}           % Plazo máximo para notificar a proveedor sobre resultados de auditoría.
\newcommand{\PlazoPlanDeAccionAuditProveedor}{2 semanas}    % Plazo maximo para recibir plan de acción de parte del proveedor tras auditoría

%%%%%%%%%%%%%%%% PSA-2.15 | Almacenamiento de registros %%%%%%%%%%%%%%%%
\newcommand{\VigenciaAlmacRegistros}{2 años}
\newcommand{\VigenciaAlmacRegistrosElec}{3 años}

%%%%%%%%%%%%%%%% GLOBAL | Resetear glosarios en cada sección %%%%%%%%%%%%%%%%
\usepackage{etoolbox}
\preto\section{\glsresetall}

\newcommand{\AC}{aseguramiento de calidad}
\newcommand{\G}{gerencia}
\newcommand{\AD}{alta dirección}
\newcommand{\OP}{operaciones}
\newcommand{\Emb}{embarques}
\newcommand{\MC}{mesa de control}     % Definición de variables estáticas
                                        % i.e. Tiempo máximo de cámras de refrigeración sin energía



%%%%%%%%%%%%%%%%%%%%%%%%%%% ESPECIALES %%%%%%%%%%%%%%%%%%%%%%%%%%%%%%%%%%
% \RequirePackage{expl3}    % \usepackage cannot be used before \documentclass
% \ExplSyntaxOn             % Switch on expl3 syntax
% \pdf_version_gset:n{1.7}  % Use provided expl3 function
% \ExplSyntaxOff            % Switch off expl3 syntax

%%%%%%%%%%%%%%%%%%%%%%%%%%% Tipografía %%%%%%%%%%%%%%%%%%%%%%%%%%%%%%%%%%

\usepackage[sfdefault,t,lf]{carlito}    % Alternativa OpenSource a Calibri, definir sans serif como predeterminado, numerales tabulares
% \usepackage[default,lnum]{opensans}
% \usepackage{CormorantGaramond}

% \usepackage{lucidabr}

% \usepackage{fontspec}

\fontsize{8.8pt}{12pt}\selectfont
\renewcommand{\baselinestretch}{1.2}

% \linespread{1.3}
% \DeclareUnicodeCharacter{2212}{-}

\usepackage{polyglossia}
	\setdefaultlanguage[spanishoperators=all]{spanish}  % Definir la localización del documento como es-MX (para que cuadro sea tabla, etc.)
	% \setdefaultlanguage{mexican}
	\renewcommand{\listtablename}{Índice de tablas}

\usepackage{changelog}  % Control de versiones
\usepackage{float}
\usepackage{pdflscape}
\usepackage{pifont}
\newcommand{\cmark}{\ding{51}}%
\newcommand{\xmark}{\ding{55}}%

% \usepackage[%
% toc=true,
% postpunc={;},
% acronyms,
% record={alsoindex}]{glossaries-extra}    % Manager de glosarios, usar xindy para el indice alfabético
% % \usepackage{glossary-mcols}
% % \renewcommand{\glsmcols}{2}
% \setglossarystyle{long3col}

% 	\makeglossaries{}
% 	% \setabbreviationstyle[acronym]{short-long}  % Estilo de acrónimos
% 	\setabbreviationstyle[acronym]{long-short}
% 	% \setglossarystyle{mcolindex}
% 	\loadglsentries[main]{../glosario.bib}  % Archvio que LuaLaTeX toma para las definiciones

% \usepackage[record]{glossaries-extra}
% 	\GlsXtrLoadResources[%
% 	src={../glosario.bib},% data in entries.bib
% 	sort={es-ES},% sort according to 'en-GB' locale
% 	]

\usepackage[%
record={alsoindex},			% create group field and other stuff
postpunc={dot},	% full stop after description
toc=true,
style=indexgroup,
acronyms,
xindy
]{glossaries-extra}

\setabbreviationstyle[acronym]{long-short}
\setglossarystyle{long3col}

\GlsXtrLoadResources[
src={../glosario.bib},	% bib files
selection={all},			% select all entries
sort={es-MX}	
]

\renewcommand{\glossarypreamble}{\fontsize{10pt}{12pt}\selectfont}

\usepackage[xindy]{imakeidx}    % Para hacer indices
	\makeindex[columns=2]   	% Estilo del indice alfabético

% \newcounter{note}[section]  % Nuevo contador para notas
% \renewcommand{\thenote}{\thesection.\Alph{note}}    % Definir estilo de numeración de las notas
% \usepackage[framemethod=TikZ]{mdframed}             % Motor gráfico para notas (TikZ)
% \newenvironment{note}[1][]{%                        % Nuevo entorno de Nota (note) 
%     \refstepcounter{note}
%     \begin{mdframed}[%                                  
%         frametitle={Nota \thenote\ #1},                 % Titulo de las notas
%         skipabove=\baselineskip{} plus 2pt minus 1pt,     % Espaciado superior
%         skipbelow=\baselineskip{} plus 2pt minus 1pt,     % Espaciado inferior
%         linewidth=0.5pt,                                % Tamaño de linea del recuadro
%         frametitlerule=true,                            % incluir titulo
%         frametitlebackgroundcolor=gray!30               % Color del recuadro
%     ]%
% }{%
%     \end{mdframed}
% }

% \newcounter{contact}[section]                           % Nuevo contador para contactos
% \renewcommand{\thecontact}{\thesection.\Alph{contact}}  % Definir que la numeración sea Sección.Alfabético

% \newenvironment{contact}[1][]{%
%     \refstepcounter{contact}
%     \begin{mdframed}[%
%         frametitle={Contacto \thecontact\ #1},          % Titulo del contacto
%         skipabove=\baselineskip{}plus 2pt minus 1pt,     % Espaciado superior
%         skipbelow=\baselineskip{}plus 2pt minus 1pt,     % Espaciado inferior
%         linewidth=0.5pt,                                % Tamaño de linea del recuadro
%         frametitlerule=true,                            % poner titulo
%         frametitlebackgroundcolor=yellow!30             % Color del recuadro
%     ]%
% }{%
%     \end{mdframed}
% }

\usepackage[multiple]{footmisc}     % Cambios en la presentación de pies de página
\usepackage{microtype}              % Ajustes de kerning y otros 
\usepackage{graphicx}               % Mejora al motor de imagenes
\usepackage{hyperref}               % Referencias cruzadas en PDF
\usepackage{cleveref}               % Referencias cruzadas en texto
%%%%%% REDEFINICIONES DE REFERENCIAS CRUZADAS %%%%%%%%%
    \crefname{note}{nota}{notas}    
	\Crefname{note}{Nota}{Notas}
	\crefname{contact}{contacto}{contactos}
	\Crefname{contact}{Contacto}{Contactos}
	\crefname{paragraph}{parrafo}{parrafos}
	\Crefname{paragraph}{Parrafo}{Parrafos}
%%%%%%%%%%%%%%%%%%%%%%%%%%%%%%%%%%%%%%%%%%%%%%%%%%%%%%%
\newcommand{\email}[1]{\href{mailto:#1}{#1}}    % Comando para hacer que se puedan clickear emails
\newcommand{\cels}[1]{\qty{#1}{\celsius}}       % Escrir celsius rapido

\usepackage{lastpage}   % Definir variable para el numero total de páginas
\usepackage{siunitx}    % Unidades internacionales y alineamiento de decimales en tablas
	\sisetup{detect-all}    % Detectar automaticamente que la tipografía es sans serif

\usepackage{color}  % definir colores nuevos

\usepackage{tabularray} % Creación de tablas
	\UseTblrLibrary{booktabs,siunitx}   % librerias parea tablas correctas, unidades internacionales
	\definecolor{Gallery}{rgb}{0.929,0.929,0.929}   % definición del color "Gallery"
	\DefTblrTemplate{contfoot-text}{normal}{\textit{Continúa en la siguiente página}}   % Continuación de tabla
	\SetTblrTemplate{contfoot-text}{normal} % Continuación de tabla normal (footer)
	\DefTblrTemplate{conthead-text}{normal}{\textit{(Continuación)}}    % Continuación de tabla
	\SetTblrTemplate{conthead-text}{normal} % Continuación de tabla (header)

    % \usepackage[margin=1cm,bmargin=2cm,tmargin=5cm,headheight=110pt]{geometry}   % Dispoción espacial del documento
% \newgeometry{includeheadfoot, heightrounded, left=1cm, right=1cm, top=1cm, bottom=4cm, headheight=6em, footskip=1.5em}    
%     \savegeometry{L1}

% \newgeometry{includeheadfoot, heightrounded, left=1cm, right=1cm, top=1cm, bottom=4.5cm, headheight=7em, footskip=1.5em, showcrop, showframe}
%     \savegeometry{L2}

\usepackage[%
includeheadfoot,
heightrounded,
left=1cm,
right=1cm,
top=1cm,
bottom=1cm,
headheight=110pt,
footskip=1.5em]{geometry}    

\usepackage[export]{adjustbox}

\usepackage{enumitem}
	\setlist[1]{labelindent=\parindent} % Indentación como parrafo para listas
	\setlist[description]{style=nextline}
	\setlist[itemize]{leftmargin=*} 	% Eliminar margen en izq
	\setlist[itemize,1]{label=---}  	% Viñeta como em-dash
	\setlist[itemize,2]{label=--}
	\setlist{noitemsep} 				% Que no se separen los item
	% [label={\theparagraph--\arabic*}] % Numerarar parrafo-numero

	\newlist{ListaSec}{enumerate}{3}
	\setlist[ListaSec]{label* = \arabic*., name* = \arabic{section}.\arabic{subsection}}


%%%%%%%%%%%%%%%% DEFINICIÓN DE CABECERA %%%%%%%%%%%%%%%%
\usepackage{fancyhdr}   % Usar fancy headers y footers
	\renewcommand{\headrulewidth}{0pt}  % Grosor de regla del header
	\fancyhead[CE,CO,LE,LO,RE,RO]{} % Resetear headers    
    \fancyfoot[CE,CO,LE,LO,RE,RO]{} % Resetear footers

%%%%%%%%%%%%%%%%%%%%%% VARIABLES DE CADA DOCUMENTO DENTRO DE UN PROGRAMA %%%%%%%%%%%%%%%%%%%%%
	\newcommand{\Codigo}{}                  % Codigo del documento
	\newcommand{\FechaPub}{}                % Fecha de publicación
	\newcommand{\Titulo}{}                  % Titulo del documento
	\newcommand{\MayorVer}{}                % Versión mayor
	\newcommand{\MenorVer}{}                % Verisón menor
	\newcommand{\Edit}{\MayorVer.\MenorVer} % Definición de la edición

    \newcommand{\Elaboro}{}                 % Nombre de quien lo elaboró
    \newcommand{\Reviso}{}                  % Nombre de quien revisó el documento
    \newcommand{\Autorizo}{}                % Nombre de quien aprobó el documento

%%%%%%%%%%%%%%%%%% Formato por defecto para informaicón documentada general %%%%%%%%%%%%%%%%%%
\fancypagestyle{formato-PS}{%
	\fancyhead{}
	\fancyfoot{}
	\fancyhead[RE]{%
		\begin{tblr}{%	
			colspec = {X[c,m]X[c,m]},    % Especificación de columnas: Anchura estática, centrado en medio
			hlines,
			vlines,
			width = 0.35\linewidth,}%
			{\textbf{Código:}\\ \Codigo} 				 & {\textbf{Edición:}\\ \MayorVer} 			\\ 	% Logotipo (2 columnas), red de frios (3 columnas), código del documento
			{\textbf{Fecha de publicación:}\\ \FechaPub} & {Página~\thepage~de~\pageref{LastPage}} 		% Titulo (4 columnas), Fecha, edición
		\end{tblr}
	}%

	\fancyhead[LO]{%
		\begin{tblr}{%	
			colspec = {X[c,m]X[c,m]},   % Especificación de columnas: Anchura estática, centrado en medio
			hlines,
			vlines,
			width = 0.35\linewidth,}%
			{\textbf{Código:}\\ \Codigo} 					& {\textbf{Edición:}\\ \MayorVer} 			\\	% Logotipo (2 columnas), red de frios (3 columnas), código del documento
			{\textbf{Fecha de publicación:}\\\FechaPub}	& {Página~\thepage~de~\pageref{LastPage}} 		% Titulo (4 columnas), Fecha, edición
		\end{tblr}
	}%

	\fancyfoot[C]{Página~\thepage~de~\pageref{LastPage}}   % Para que muestre Página # de #
	\fancyfoot[L]{\textbf{CONFIDENCIAL}}                    % CONFIDENCIAL
	\fancyfoot[R]{\Codigo~v.\Edit}                          % Estructura del código CODIGO v.EDICION
}%

\fancypagestyle{formato-PI}{%
	\fancyhead{}
	\fancyfoot{}
    \fancyhead[C]{%
	    \begin{tblr}{%	
            colspec = {X[c,m]X[c,m]X[c,m]X[c,m]X[c,m]X[c,m]},    % Especificación de columnas: Anchura estática, centrado en medio
            hlines,
            vlines,
			width = \linewidth,}%
			\SetCell[c=2]{} \includegraphics[height=3em, valign=c]{./RDF_Logo.png}	&	& \SetCell[c=3]{} \large{Red de fríos S.A. de C.V.}	&	&												& {\textbf{Código:}\\ \Codigo}	\\ % Logotipo (2 columnas), red de frios (3 columnas), código del documento
			\SetCell[c=4]{} {\textbf{Titulo:}\\ \Titulo} 							&	&													&	& {\textbf{Fecha de publicación:}\\ \FechaPub}	& {\textbf{Edición:}\\ \Edit}	\\ % Titulo (4 columnas), Fecha, edición
            \end{tblr}
	}
    \fancyfoot[C]{Página~\thepage~de~\pageref{LastPage}}   % Para que muestre Página # de #
    \fancyfoot[L]{\textbf{CONFIDENCIAL}}                    % CONFIDENCIAL
    \fancyfoot[R]{\Codigo~v.\Edit}                          % Estructura del código CODIGO v.EDICION
}%

%%%%%%%%%%%%%%%%%%%%%%%%%%%%%%%%%%%%%%%%%%%%%%%%%%%%%%%%%%%%%%%%%%%%%%%%%%%%%%%%%%%%%%%%%%%%%%%%%%%%%%%%%%%%%%%%%%
			
\setcounter{tocdepth}{4}            % Qué tanto mostrará la tabla de contenidos
\setcounter{secnumdepth}{6}         % Que tantas secciones numerará LaTeX (Se ajustó hasta subparrafo)

\newcommand{\defglo}[1]{\item[\glsname{#1}] \glsdesc{#1}} % Define un comando para dar un item en un entorno de descripción para una palabra definida en el glosario

% \includeonly{src/HyS-3.25.tex}

\begin{document}
\pagestyle{formato-PS}
\thispagestyle{formato-PI}
\renewcommand{\MayorVer}{1}
\renewcommand{\MenorVer}{0}
\renewcommand{\Codigo}{HyS-ToC}
\renewcommand{\FechaPub}{2023--01}
%\renewcommand{\Edit}{01}
\renewcommand{\Titulo}{Programa de Buenas Prácticas de Distribución --- Índice}

\tableofcontents
\thispagestyle{formato-PI}
\renewcommand{\MayorVer}{2}
\renewcommand{\MenorVer}{0}
%\renewcommand{\Edit}{\MayorVer.\MenorVer}
\renewcommand{\Titulo}{Limpieza de cámaras de refrigeración}
\renewcommand{\Codigo}{HYS-1-IT}
\renewcommand{\FechaPub}{2023--01}
% \renewcommand{\Titulo}{Politica de calidad}

\section{\Titulo}\index{Limpieza de!cámara!refrigeración, de}

\subsection{Objetivo}

\begin{itemize}
    \item \textbf{Establecer} el procedimiento de limpieza que se debe de realizar en una cámara de refrigeración y congelación, buscando con ello eliminar o minimizar la presencia de peligros físicos, químicos o biológicos que pudieran generar un riesgo de contaminación en el producto.
\end{itemize}


\subsection{Alcance}
Este procedimiento es aplicable a todas las cámaras en temperaturas de refrigeración y congelación de \gls{RDF}

\subsection{Terminología y definiciones}
\begin{description}
    \defglo{limpieza}
    \defglo{sanitizacion}
\end{description}

% \subsection{Documentos y/o normas relacionadas} %TODO

\subsection{Procedimiento}
\subsubsection{Materiales}
\begin{itemize}
    \item Preparación de KOOL KLEAN a la concentración indicada por el proveedor (Aplicación directa / NO diluya).
          \begin{itemize}
              \item Cubeta de \qty{20}{\liter}
              \item Cepillo para piso color rojo
              \item Cepillos manuales color rojo
              \item Recogedor de plástico
              \item Jalador color rojo
              \item Trapo de microfibra antibacterial.
          \end{itemize}
\end{itemize}

\subsubsection{Instrucciones}

\paragraph{Limpieza preoperativa}

\begin{enumerate}
    \item Retire el producto del área donde se realizara la limpieza para evitar una posible contaminación durante el proceso de limpieza.
    \item Retire el material o suciedad visible como pudiera ser emplaye, madera, cartón, papel entre otros hasta que el área quede libre de estos, deposite los desechos en el contenedor de basura del área.
    \item Utilizar detergente \emph{KOOL KLEANE} al \qty{100}{\percent} (no diluya) para la limpieza de techos, muros, \textit{racks} y pisos.
    \item La limpieza del área se debe realizar en el siguiente orden: techos, muros, \textit{racks} y pisos.
    \item Primero, realice la limpieza de techos y paredes. Aplique el detergente \emph{KOOL KLEANE} directamente de manera superficial con un trapo de microfibra, deje el producto aplicado durante \qtyrange{10}{20}{\minute} aproximadamente para que genere su acción
    \item Con otro trapo de microfibra retire el excedente de detergente, realice esta maniobra hasta finalizar el área seleccionada para la limpieza.
    \item Elimine del trapo de microfibra el excedente de detergente y continúe eliminando el excedente hasta finalizar de retirar el detergente del área.
    \item Continúe con la limpieza de los \textit{racks}, aplique el detergente \emph{KOOL KLEANE} directamente de manera superficial con un trapo de microfibra, deje el producto aplicado durante \qtyrange{10}{20}{\minute} aproximadamente para que genere su acción
    \item Con otro trapo de microfibra retire el excedente de detergente, realice esta maniobra hasta finalizar el área seleccionada para la limpieza.
    \item Elimine del trapo de microfibra el excedente de detergente y continúe eliminando el excedente hasta finalizar de retirar el detergente del área.
    \item Finalice con la limpieza de los pisos y zoclos, aplique el detergente \emph{KOOL KLEANE} en el piso con ayuda de la escoba para piso, deje el producto aplicado durante \qtyrange{10}{20}{\minute} aproximadamente para que genere su acción.
    \item Talle con la escoba el piso y zoclos.
    \item Con el jalador para pisos retire el excedente de detergente hasta que el área quede libre del mismo.
    \item Una vez que se encuentren libres de detergente las áreas limpiadas
    \item Registre en el \emph{Programa Maestro de Limpieza y Sanitización} (\RAC. la actividad realizada.
    \item Finalmente, notifique al Supervisor de Calidad para la verificación y liberación de la tarea.
    \item La tarea es liberada una vez que el área se encuentran libre de polvo, suciedad, residuos de materiales extraños, residuos de jabón, y seca. En caso contrario la tarea se deberá volver a realizar.
\end{enumerate}

\paragraph{Limpieza operativa}

\begin{enumerate}
    \item Con la escoba y recogedor retire el material o suciedad visible como pudiera ser emplaye, madera, cartón, papel entre otros hasta que el piso del área quede libre de estos.
    \item Retire de manera manual el emplaye, madera, cartón, papel que pudieran estar acumulados entre los \textit{racks}
    \item Deposite los desechos en el contenedor de basura del área.
\end{enumerate}

\subsection{Responsables de la actividad}

\begin{itemize}
    \item \textbf{Ejecutado} por personal de limpieza
    \item \textbf{Monitoreado} por personal de mantenimiento
    \item \textbf{Verificado} por personal de Calidad.
\end{itemize}

\subsection{Acciones preventivas}

\begin{itemize}
    \item Se llevará a cabo una inspección diaria. Registrando el indicador y medida correspondiente
    \item Si la desviación se repite frecuentemente se dará curso de capacitación al personal de almacén y limpieza para que realicen eficientemente su trabajo.
    \item Si después de haber capacitado al personal de almacén y limpieza se siguen presentando desviaciones por causas injustificadas, se tomaran acciones más enérgicas con el personal por incumplimiento con sus deberes.
\end{itemize}

\subsection{Acciones correctivas}

\begin{itemize}
    \item La tarea es liberada una vez que el área se encuentran libre de polvo, suciedad, residuos de materiales extraños, residuos de jabón, y seca. En caso contrario la tarea se deberá volver a realizar como se indica en el procedimiento.
    \item En caso de no conformidad reportar en \RAC.
\end{itemize}

\subsection{Frecuencia}
\begin{itemize}
    \item[Superficial:] diario de lunes a viernes de 8:00 a 9:00, sábados de 9:00 a 10:00.
    \item[Profunda:] Anualmente, de acuerdo con el programa y disponibilidad de las cámaras. 
\end{itemize}

\begin{changelog}[simple, sectioncmd=\subsection*,label=changelog-3.1]
	\begin{version}[v=2.0, date=2023--01, author=Pablo E. Alanis]
		\item Cambio de formato;
		\item Cambios en la serialización de versiones;
		\item Correcciones ortográficas y de estilo.
	\end{version}

	\begin{version}[v=1.5, date=2022--05, author=Alonso M.]
		\item cambio de fecha;
	\end{version}

	\shortversion{v=1.4, date=2021--05, changes=No hubo cambios}
\end{changelog}

\thispagestyle{formato-PI}
\renewcommand{\MayorVer}{2}
\renewcommand{\MenorVer}{0}
%\renewcommand{\Edit}{\MayorVer.\MenorVer}
\renewcommand{\Titulo}{Limpieza de cámara de secos}
\renewcommand{\Codigo}{HYS-2-IT}
\renewcommand{\FechaPub}{2023--01}
% \renewcommand{\Titulo}{Politica de calidad}

\section{\Titulo}\index{Limpieza de!cámara!secos, de}

% \section{}

\subsection{Objetivo}

\begin{itemize}
	\item \textbf{Establecer} el procedimiento de limpieza que se debe de realizar en el almacén de secos, buscando con ello eliminar o minimizar la presencia de peligros físicos, químicos o biológicos que pudieran generar un riesgo de contaminación en el producto.
\end{itemize}

\subsection{Alcance}

Este procedimiento es aplicable a todos los almacenes de secos de \gls{RDF}.

\subsection{Terminología y definiciones}

\begin{description}
	\defglo{limpieza}
	\defglo{sanitizacion}
\end{description}

% \subsection{Documentos y/o normas relacionadas}

\subsection{Procedimiento}
\subsubsection{Materiales}
\begin{itemize}
	\item Preparación de Liquid-K a la concentración indicada por el proveedor (\qtyrange{2.6}{3.0}{\milli\liter\per\liter})
	\item Cubeta
	\item Cepillo para piso color rojo
	\item Cepillos manuales color rojo
	\item Recogedor de plástico
	\item Jalador color rojo
	\item Trapo de microfibra antibacteriana.
\end{itemize}

\subsubsection{Instrucciones}
\paragraph{Limpieza}
\begin{enumerate}
	\item Retire el producto del área donde se realizara la limpieza para evitar una posible contaminación durante el proceso de limpieza.
	\item Retire el material o suciedad visible como pudiera ser emplaye, madera, cartón, papel entre otros hasta que el área quede libre de estos, deposite los desechos en el contenedor de basura del área
	\item Utilizar detergente \textit{Liquid K} (\qtyrange{1.3}{2.6}{\milli\liter\per\liter}) para la limpieza de techos, muros, \textit{racks} y pisos.
	\item La limpieza del área se debe realizar en el siguiente orden: techos, muros, \textit{racks} y pisos
	\item Primero realice la limpieza de techos y paredes, aplique el detergente \textit{Liquid K} directamente de manera superficial con un trapo de microfibra, deje el producto aplicado durante \qtyrange{10}{20}{\minute} aproximadamente para que genere su acción.
	\item Con otro trapo de microfibra retire el excedente de detergente, realice esta maniobra hasta finalizar el área seleccionada para la limpieza.
	\item Elimine del trapo de microfibra el excedente de detergente y continúe eliminando el excedente hasta finalizar de retirar el detergente del área.
	\item Continúe con la limpieza de los \textit{racks}, aplique el detergente \textit{Liquid K} de manera superficial con un trapo de microfibra, deje el producto aplicado durante \qtyrange{10}{20}{\minute} aproximadamente para que genere su acción
	\item Con otro trapo de microfibra retire el excedente de detergente, realice esta maniobra hasta finalizar el área seleccionada para la limpieza.
	\item Elimine del trapo de microfibra el excedente de detergente y continúe eliminando el excedente hasta finalizar de retirar el detergente del área
	\item Finalice con la limpieza de los pisos y zoclos, aplique el detergente \textit{Liquid K} en el piso con ayuda de la escoba para piso, deje el producto aplicado durante \qtyrange{10}{20}{\minute} aproximadamente para que genere su acción.
	\item Talle con la escoba el piso y zoclos.
	\item Con el jalador para pisos retire el excedente de detergente hasta que el área quede libre del mismo.
	\item Una vez que se encuentren libres de detergente las áreas limpiadas
	\item Registre en el Programa Maestro de Limpieza y Sanitización (\RAC. la actividad realizada.
	\item Finalmente, notifique al Supervisor de Calidad para la verificación y liberación de la tarea.
	\item La tarea es liberada una vez que el área se encuentran libre de polvo, suciedad, residuos de materiales extraños, residuos de jabón, y seca. En caso contrario la tarea se deberá volver a realizar.
\end{enumerate}

\subsection{Responsables de la actividad}
\begin{itemize}
	\item \textbf{Ejecutado} por personal de limpieza
	\item \textbf{Monitoreado} por personal de mantenimiento
	\item \textbf{Verificado} por personal de Calidad.
\end{itemize}

\subsection{Acciones preventivas}
\begin{enumerate}
	\item La tarea es liberada una vez que el área se encuentran libre de polvo, suciedad, residuos de materiales extraños, residuos de jabón, y seca. En caso contrario la tarea se deberá volver a realizar como se indica en el procedimiento.
	\item En caso de no conformidad reportar en \RAC.
\end{enumerate}

\subsection{Frecuencia}
\begin{itemize}
	\item[\textbf{Superficial:}] diario de lunes a viernes de 8:00 a 9:00, sábados de 9:00 a 10:00.
	\item[\textbf{Profunda:}] Anualmente, de acuerdo con el programa y disponibilidad de las cámaras.
\end{itemize}

\begin{changelog}[simple, sectioncmd=\subsection*,label=changelog-3.2]
	\begin{version}[v=2.0, date=2023--01, author=Pablo E. Alanis]
		\item Cambio de formato;
		\item Cambios en la serialización de versiones;
		\item Correcciones ortográficas y de estilo.
	\end{version}

	\begin{version}[v=1.5, date=2022--05, author=Alonso M.]
		\item cambio de fecha;
	\end{version}

	\shortversion{v=1.4, date=2021--05, changes=No hubo cambios}
\end{changelog}
\thispagestyle{formato-PI}
\renewcommand{\MayorVer}{2}
\renewcommand{\MenorVer}{0}
%\renewcommand{\Edit}{\MayorVer.\MenorVer}
\renewcommand{\Titulo}{Limpieza de andén y cámaras desocupadas}
\renewcommand{\Codigo}{HYS-3-IT}
\renewcommand{\FechaPub}{2023--01}

\section{\Titulo}\index{Limpieza!andén, de}\index{Limpieza!cámaras desocupadas, de}
% \section{Limpieza de andén y cámaras desocupadas}

\subsection{Objetivo}
\begin{itemize}
	\item \textbf{Establecer} el procedimiento de limpieza que se debe de realizar en andén y cámara fuera de operación, buscando con ello eliminar o minimizar la presencia de peligros físicos, químicos o biológicos que pudieran generar un riesgo de contaminación en el producto.
\end{itemize}

\subsection{Alcance}
Este procedimiento es aplicable al andén y cámaras fuera de operación de \gls{RDF}.

\subsection{Terminología y definiciones}

\begin{description}
\defglo{limpieza}
\defglo{sanitizacion}
\end{description}

% \subsection{Documentos y/o normas relacionadas} %TODO

\subsection{Procedimiento}
\subsubsection{Materiales}

\begin{itemize}
	\item Preparación de \textit{Liquid K} a la concentración indicada por el proveedor (\qtyrange{1.3}{2.6}{\milli\liter\per\liter}).
	\item Cubeta
	\item Cepillo para piso color rojo
	\item Cepillos manuales color rojo
	\item Recogedor de plástico
	\item Jalador color rojo
Trapo de microfibra antibacterial.
\end{itemize}

\subsubsection{Instrucciones}
\paragraph{Preoperativo}
\begin{enumerate}
	\item Retire el producto del área donde se realizara la limpieza para evitar una posible contaminación durante el proceso de limpieza.
	\item Retire el material o suciedad visible como pudiera ser emplaye, madera, cartón, papel, entre otros hasta que el área quede libre de estos, deposite los desechos en el contenedor de basura del área.
	\item Con ayuda de la hidrolavadora pre enjuague los techos, paredes, protecciones, puertas, \textit{racks}, cortinas y pisos.
	\item En una cubeta preparar la solución de detergente \textit{Liquid K} (\qtyrange{1.3}{2.6}{\milli\liter\per\liter}).
	\item La limpieza del área se debe realizar en el siguiente orden: techos, muros, protecciones, puertas \textit{racks}, cortinas y pisos.\label{3.3:pt5}
	\item Sumerja el cepillo en el detergente \textit{Liquid K} y talle el área, realice esta operación las veces que sean necesarias hasta cubrir toda el área seleccionada, inicie el tallado de arriba hacia abajo según indica el orden en \cref{3.3:pt5}.
	\item Para el caso del piso, talle con el cepillo destinado para esta área.
	\item Con el uso de la hidrolavadora, realice el enjuague del área iniciando de la parte más alta y terminado con los pisos.
	\item Con el jalador rastrille el excedente de agua de los pisos y retire el agua con la ayuda de la aspiradora.
	\item Registre en el Programa Maestro de Limpieza y Sanitización (\RAC. la actividad realizada.
	\item Finalmente, notifique al Supervisor de Calidad para la verificación y liberación de la tarea.
	\item La tarea es liberada una vez que el área se encuentran libre de polvo, suciedad, residuos de materiales extraños, residuos de jabón, y seca. En caso contrario la tarea se deberá volver a realizar
	\item Con la escoba y recogedor retire el material o suciedad visible como pudiera ser emplaye, madera, cartón, papel entre otros hasta que el piso del área quede libre de estos;
	\item Retire de manera manual el emplaye, madera, cartón, papel que pudieran estar acumulados entre los \textit{racks.}
	\item Deposite los desechos en el contenedor de basura del área.
\end{enumerate}

\subsection{Responsables de la actividad}

\begin{itemize}
	\item \textbf{Ejecutado} por personal de limpieza
	\item \textbf{Monitoreado} por personal de mantenimiento
	\item \textbf{Verificado} por personal de Calidad.
\end{itemize}

\subsection{Acciones preventivas} %TODO: Hay que especificar que es lo que se verifica a diario.

\begin{itemize}
	\item Se llevará a cabo una inspección diaria. Registrando el indicador y medida correspondiente
	\item Si la desviación se repite frecuentemente se dará curso de capacitación al personal de almacén y limpieza para que realicen eficientemente su trabajo.
	\item Si después de haber capacitado al personal de almacén y limpieza se siguen presentando desviaciones por causas injustificadas, se tomaran acciones más enérgicas con el personal por incumplimiento con sus deberes.
\end{itemize}

\subsection{Acciones correctivas}

\begin{itemize}
	\item La tarea es liberada una vez que el área se encuentran libre de polvo, suciedad, residuos de materiales extraños, residuos de jabón, y seca. En caso contrario la tarea se deberá volver a realizar como se indica en el procedimiento. 
	\item En caso de no conformidad reportar en \RAC.
\end{itemize}

\subsection{Frecuencia}

\begin{itemize}
	\item[\textbf{Superficial:}]	Diaria, de lunes a viernes de 9:00 a 10:00 y sábados de 10:00 a 11:00;
	\item[\textbf{Profunda:}]		Trimestral según el programa general de limpieza.\footnote{En caso de que no haya disponibilidad para la limpieza profunda del andén en el intervalo establecido por el programa de limpieza, esta se puede reprogramar en un máximo de 2 meses.}
\end{itemize}

\begin{changelog}[simple, sectioncmd=\subsection*,label=changelog-\thesection-\MayorVer.\MenorVer]
	\begin{version}[v=\MayorVer.\MenorVer, date=2023--01, author=Pablo E. Alanis]
		\item Cambio de formato;
		\item Cambios en la serialización de versiones;
		\item Correcciones ortográficas y de estilo.
	\end{version}

	\begin{version}[v=1.5, date=2022--05, author=Alonso M.]
		\item cambio de fecha;
	\end{version}

	\shortversion{v=1.4, date=2021--05, changes=No hubo cambios}
\end{changelog}
\thispagestyle{formato-PI}
\renewcommand{\MayorVer}{2}
\renewcommand{\MenorVer}{0}
%\renewcommand{\Edit}{\MayorVer.\MenorVer}
\renewcommand{\Codigo}{HYS-4-IT}
\renewcommand{\FechaPub}{2023--01}
\renewcommand{\Titulo}{Limpieza de rampas mecánicas}

\section{\Titulo}\index{Limpieza de!rampas mecánicas, de}

% \section{Limpieza de rampas mecánicas}

\subsection{Objetivo}

\begin{itemize}
	\item \textbf{Establecer} el procedimiento de limpieza que se debe de realizar en a las rampas mecánicas para minimizar la presencia de peligros físicos, químicos o biológicos que pudieran generar un riesgo al producto.
\end{itemize}

\subsection{Alcance}
Todas las rampas mecánicas ubicadas en el almacén.

\subsection{Terminología y definiciones}

\begin{description}
	\defglo{limpieza}
	\defglo{sanitizacion}
\end{description}

%\subsection{Documentos y/o normas relacionadas}

%N/A

\subsection{Procedimiento}
\subsubsection{Materiales}
\begin{itemize}
	\item Cepillo para piso color rojo
	\item Cepillos manuales color rojo
	\item Recogedor de plástico
\end{itemize}

\subsubsection{Instrucciones}
\paragraph{Preoperativo}

\begin{enumerate}
	\item Poner el cono de aviso, levantar rampa y poner el seguro de la rampa.
	\item Retire con ayuda de escoba todos los residuos que se encuentren en la superficie y en el interior de la rampa.
	\item Posteriormente se procede a la recolección de los residuos con recogedor.
	\item Depositar los desechos en el contenedor para basura.
	\item Retire con ayuda del cepillo todos los residuos que se encuentran en la superficie y en el interior de la rampa.
	\item Posteriormente se procede a la recolección de los residuos con recogedor.
	\item Depositar los desechos en el contenedor para basura.
\end{enumerate}

\paragraph{Operativo}
\begin{enumerate}
	\item Retirar con ayuda de cepillo para piso todos los residuos (plásticos, pedazos de madera, etc.).
	\item Retire residuos de los cepillos de la rampa
	\item Posteriormente se procede a la recolección de los residuos con recogedor.
	\item Depositar los desechos en el contenedor para basura.
\end{enumerate}

\subsection{Responsables de la actividad}
\begin{itemize}
	\item \textbf{Ejecutado} por personal de limpieza
	\item \textbf{Monitoreado} por personal de mantenimiento
	\item \textbf{Verificado} por personal de Calidad.
\end{itemize}

\subsection{Acciones preventivas}

\begin{itemize}
	\item Se llevará a cabo una inspección diaria. Registrando el indicador y medida correspondiente
	\item Si la desviación se repite frecuentemente se dará curso de capacitación al personal de limpieza para que realicen eficientemente su trabajo.
	\item Si después de haber capacitado al personal de limpieza se siguen presentando desviaciones por causas injustificadas, se tomaran acciones más enérgicas con el personal por incumplimiento con sus deberes.
\end{itemize}

\subsection{Acciones correctivas}

\begin{itemize}
	\item La tarea es liberada una vez que el área se encuentran libre de polvo, suciedad, residuos de materiales extraños, residuos de jabón, y seca. En caso contrario la tarea se deberá volver a realizar como se indica en el procedimiento.
	\item En caso de no conformidad reportar en \RAC.
\end{itemize}

\subsection{Frecuencia}

\begin{itemize}
	\item[\textbf{Superficial}] Diaria de lunes a viernes de 10:30 a 11:00 y sábados de 11:30 a 12:00;
	\item[\textbf{Profunda:}] Anualmente, según el programa y la disponibilidad de espacio.
\end{itemize}

\begin{changelog}[simple, sectioncmd=\subsection*,label=changelog-3.7]
	\begin{version}[v=2.0, date=2023--01, author=Pablo E. Alanis]
		\item Cambio de formato;
		\item Cambios en la serialización de versiones;
		\item Correcciones ortográficas y de estilo.
	\end{version}

	\begin{version}[v=1.7, date=2022--05, author=Alonso M.]
		\item cambio de fecha;
	\end{version}

	\shortversion{v=1.6, date=2021--05, changes=No hubo cambios}
\end{changelog}
\thispagestyle{formato-PI}
\renewcommand{\MayorVer}{2}
\renewcommand{\MenorVer}{0}
%\renewcommand{\Edit}{\MayorVer.\MenorVer}
\renewcommand{\Titulo}{Limpieza de montacargas y patines}
\renewcommand{\Codigo}{HYS-5-IT}
\renewcommand{\FechaPub}{2023--01}

\section{\Titulo}\index{Limpieza!montacargas, de}\index{Limpieza!patines, de}
% \section{Limpieza de montacargas y patines}

\subsection{Objetivo}
Establecer el procedimiento de limpieza que se debe de realizar a los montacargas y patines de carga para su mantenimiento y minimizar riesgos de contaminación al producto.

\subsection{Alcance}
Todos los patines y montacargas ubicados dentro del almacén.

\subsection{Terminología y definiciones}
\begin{description}
	\defglo{limpieza}
	\defglo{sanitizacion}
\end{description}

\subsection{Procedimiento}

\subsubsection{Materiales}

\begin{itemize}
	\item Preparación de \textit{Liquid K} a la concentración indicada por el proveedor \qtyrange{0.13}{0.26}{\percent} o bien \qtyrange{1.3}{2.6}{\milli\liter\per\liter}
	\item Cubeta de \qty{20}{\litre}
	\item Toalla desechable
	\item Cepillo de mano
\end{itemize}

\subsubsection{Instrucciones}

\paragraph{Limpieza de montacargas y patines diario}

\subparagraph{Preoperativo}

\begin{enumerate}
	\item Desconecte equipos para evitar que se encuentren energizados.
	\item Prepare el detergente a la concentración recomendada por el proveedor, sumerja la toalla desechable en el detergente, elimine el exceso del mismo exprimiendo la toalla desechable limpie la estructura del equipo de carga.
	\item Elimine el excedente de humedad del equipo de carga utilizando toalla desechable.
	\item Limpiar el área después de terminar el trabajo.
\end{enumerate}

\subsection{Responsables de la actividad}

\begin{itemize}
	\item \textbf{Ejecutado} por personal de operación montacargas
	\item \textbf{Monitoreado} por personal de mantenimiento
	\item \textbf{Verificado} por personal de Calidad.
\end{itemize}

\subsection{Acciones preventivas}

\begin{itemize}
	\item Se llevará a cabo una inspección diaria. Registrando el indicador y medida correspondiente
	\item Si la desviación se repite frecuentemente se dará curso de capacitación al personal de almacén para que realicen eficientemente su trabajo.
	\item Si después de haber capacitado al personal de almacén se siguen presentando desviaciones por causas injustificadas, se tomaran acciones más enérgicas con el personal por incumplimiento con sus deberes.
\end{itemize}

\subsection{Acciones correctivas}

\begin{enumerate}
	\item La tarea es liberada una vez que el área se encuentran libre de polvo, suciedad, residuos de materiales extraños, residuos de jabón, y seca. En caso contrario la tarea se deberá volver a realizar como se indica en el procedimiento.
	\item En caso de no conformidad reportar en \RAC.
\end{enumerate}

\subsection{Frecuencia}

\begin{itemize}
	\item Diario al finalizar operaciones lunes a viernes 5:00 pm y sábado a la 1:00 pm
\end{itemize}

\begin{changelog}[title=Registro de cambios,simple, sectioncmd=\subsection*,label=changelog-\thesection-\MayorVer.\MenorVer]
	\begin{version}[v=\MayorVer.\MenorVer, date=2023--01, author=Pablo E. Alanis]
		\item Cambio de formato;
		\item Cambios en la serialización de versiones;
		\item Correcciones ortográficas y de estilo.
	\end{version}

	\begin{version}[v=1.7, date=2022--05, author=Alonso M.]
		\item cambio de fecha;
	\end{version}

	\shortversion{v=1.6, date=2021--05, changes=No hubo cambios}
\end{changelog}
\thispagestyle{formato-PI}
\renewcommand{\MayorVer}{2}
\renewcommand{\MenorVer}{0}
%\renewcommand{\Edit}{\MayorVer.\MenorVer}
\renewcommand{\Titulo}{Limpieza de artículos de andén}
\renewcommand{\Codigo}{HYS-6-IT}
\renewcommand{\FechaPub}{2023--01}

\section{\Titulo}\index{Limpieza de!artículos del andén}
% \section{Limpieza de artículos de andén}

\subsection{Objetivo}

Establecer un procedimiento que aseguren un mantenimiento y limpieza apropiada de los artículos de andén.

\subsection{Alcance}

Mesa de inspección y lavamanos.

\subsection{Terminología y definiciones}

\begin{description}
\defglo{limpieza}
\defglo{sanitizacion}
\end{description}

\subsection{Procedimiento}

\subsubsection{Materiales}

\begin{itemize}
	\item Preparación de Liquid K a la concentración indicada por el proveedor \qtyrange{0.13}{0.26}{\percent} o bien \qtyrange{1.3}{2.6}{\milli\liter\per\liter}
	\item Cubeta de \qty{20}{\litre}
	\item Cepillos manuales color rojo
	\item Toallas desechables
\end{itemize}

\subsubsection{Instrucciones}
\paragraph{Mesa de trabajo}
\subparagraph{Preoperativo}
Cuando la mesa y lavamanos no sea utilizada Se llevará a cabo la \emph{limpieza en seco,} que consiste en rociar desinfectante con el aspersor manual en la superficie de la mesa y lavamanos y con una toalla desechable frotar hasta secar el mismo.

\section{Mesa después de inspección}

\begin{enumerate}
	\item En caso de que se llegara a presentar derrame de residuos orgánicos (líquido y/o pedazos de algún producto cárnico), se recoge los residuos y se depositan en el contenedor.
	\item Pre enjuagar la mesa y lavamanos, se procede a aplicar el desengrasante en el área específica a la concentración indicada por el proveedor.
	\item Tallando con el cepillo iniciando de la parte superior hasta la inferior.
	\item Enjuague con agua hasta eliminar el detergente.
	\item Secar con toalla desechable.
	\item Desinfectar con la ayuda de un Atomizador rociando la Mesa
	\item Frotar la Mesa con una toalla desechable hasta que quede seca
\end{enumerate}

\subsection{Responsables de la actividad}

\begin{itemize}
	\item \textbf{Ejecutado} por personal de limpieza
	\item \textbf{Monitoreado} por personal de mantenimiento
	\item \textbf{Verificado} por personal de Calidad.
\end{itemize}

\subsection{Acciones preventivas}

\begin{itemize}
	\item Se llevará a cabo una inspección diaria. Registrando el indicador y medida correspondiente
	\item Si la desviación se repite frecuentemente se dará curso de capacitación al personal de limpieza para que realicen eficientemente su trabajo.
	\item Si después de haber capacitado al personal de limpieza se siguen presentando desviaciones por causas injustificadas, se tomaran acciones más enérgicas con el personal por incumplimiento con sus deberes.
\end{itemize}

\subsection{Acciones correctivas}

\begin{enumerate}
	\item La tarea es liberada una vez que el área se encuentran libre de polvo, suciedad, residuos de materiales extraños, residuos de jabón, y seca. En caso contrario la tarea se deberá volver a realizar como se indica en el procedimiento.
	\item En caso de no conformidad reportar en \RAC.
\end{enumerate}

\subsection{Frecuencia}

\begin{itemize}
	\item Diaria de lunes a viernes de 16:00 a 16:30 y sábados de 12:00 a 12:30.
\end{itemize}

\begin{changelog}[simple, sectioncmd=\subsection*,label=changelog-3.6]
	\begin{version}[v=2.0, date=2023--01, author=Pablo E. Alanis]
		\item Cambio de formato;
		\item Cambios en la serialización de versiones;
		\item Correcciones ortográficas y de estilo.
	\end{version}

	\begin{version}[v=1.7, date=2022--05, author=Alonso M.]
		\item cambio de fecha;
	\end{version}

	\shortversion{v=1.6, date=2021--05, changes=No hubo cambios}
\end{changelog}
\thispagestyle{formato-PI}
\renewcommand{\MayorVer}{2}
\renewcommand{\MenorVer}{0}
%\renewcommand{\Edit}{\MayorVer.\MenorVer}
\renewcommand{\Codigo}{HYS-7-IT}
\renewcommand{\FechaPub}{2023--01}
\renewcommand{\Titulo}{Limpieza de aduana sanitaria}

\section{\Titulo}\index{Limpieza!aduana sanitaria, de}

\subsection{Objetivos}
Establecer el procedimiento de limpieza que aseguren un mantenimiento y limpieza apropiada de la aduana sanitaria.

\subsection{Alcance}
Este procedimiento es aplicable a la aduana sanitaria.

\subsection{Terminología y definiciones}
\begin{description}
\defglo{limpieza}
\defglo{sanitizacion}
\end{description}

\subsection{Procedimiento}
\subsubsection{Materiales}

\begin{itemize}
	\item Preparación de Liquid K a la concentración indicada por el proveedor \qtyrange{1.3}{2.6}{\percent} o bien \qtyrange{1.3}{2.6}{\milli\liter\per\liter}.
	\item Cubeta de \qty{20}{\litre}
	\item Cepillo para piso color rojo
	\item Cepillos manuales color rojo
	\item Recogedor de plástico
	\item Jalador color rojo
	\item Toallas desechables
\end{itemize}

\subsubsection{Instrucciones}

\paragraph{Preoperativo}

\subparagraph{Limpieza de aduana sanitaria}

\begin{enumerate}
	\item Prepare detergente a la concentración indicada por el proveedor.
	\item Con ayuda de un cepillo tallar el techo, paredes, tarja y accesorios de la aduana.
	\item Con agua corriente eliminar el excedente de detergente.
	\item Talle los pisos y junta sanitaria con ayuda de cepillo.
	\item Elimine el exceso de detergente con ayuda de agua corriente.
	\item Con el jalador elimine el excedente de agua hasta dejar el área libre de humedad.
\end{enumerate}

\paragraph{Limpieza de tapete sanitario, tarja y pisos}

\subparagraph{Pisos}

\begin{enumerate}
	\item Retire los desechos con ayuda de cepillo para piso.
	\item Deposite los desechos en contenedor.
	\item Prepare detergente a la concentración indicada por el proveedor.
	\item Con ayuda de cepillo, realice limpieza de pisos.
	\item Elimine excedente de humedad.
\end{enumerate}

\subparagraph{Tapete}

\begin{enumerate}
	\item Elimine el exceso de desinfectante
	\item Con agua corriente enjuague el tapete sanitario hasta eliminar cualquier residuo
\end{enumerate}

\subparagraph{Tarja}

\begin{enumerate}
	\item Limpie la tarja, elimine cualquier desecho que se encuentre en la misma.
	\item Prepare detergente a la concentración indicada por el proveedor.
	\item Talle con cepillo de mano la tarja
	\item Enjuague con agua corriente hasta elimine el excedente de detergente.
\end{enumerate}

\subsection{Responsables de la actividad}

\begin{itemize}
	\item \textbf{Ejecutado} por personal de limpieza;
	\item \textbf{Monitoreado} por personal de mantenimiento;
	\item \textbf{Verificado} por personal de Calidad.
\end{itemize}

\subsection{Acciones preventivas}

\begin{itemize}
	\item Se llevará a cabo una inspección diaria. Registrando el indicador y medida correspondiente
	\item Si la desviación se repite frecuentemente se dará curso de capacitación al personal de limpieza para que realicen eficientemente su trabajo.
	\item Si después de haber capacitado al personal de limpieza se siguen presentando desviaciones por causas injustificadas, se tomaran acciones más enérgicas con el personal por incumplimiento con sus deberes.
\end{itemize}

\subsection{Acciones correctivas}

\begin{itemize}
	\item La tarea es liberada una vez que el área se encuentran libre de polvo, suciedad, residuos de materiales extraños, residuos de jabón, y seca. En caso contrario la tarea se deberá volver a realizar como se indica en el procedimiento.
	\item En caso de no conformidad reportar en \RAC.
\end{itemize}

\subsection{Frecuencia}

\begin{description}
	\item[Profundo] Trimestralmente;
	\item[Superficial] Lunes a sábado de 8:00 a 8:30.
\end{description}

\begin{changelog}[simple, sectioncmd=\subsection*,label=changelog-\thesection-\MayorVer.\MenorVer]
	\begin{version}[v=\MayorVer.\MenorVer, date=2023--01, author=Pablo E. Alanis]
		\item Cambio de formato;
		\item Cambios en la serialización de versiones;
		\item Correcciones ortográficas y de estilo.
	\end{version}

	\begin{version}[v=1.5, date=2022--05, author=Alonso M.]
		\item cambio de fecha;
	\end{version}

	\shortversion{v=1.4, date=2021--05, changes=No hubo cambios}
\end{changelog}
\thispagestyle{formato-PI}
\renewcommand{\MayorVer}{2}
\renewcommand{\MenorVer}{0}
%\renewcommand{\Edit}{\MayorVer.\MenorVer}
\renewcommand{\Codigo}{HYS-8-IT}
\renewcommand{\FechaPub}{2023--01}
\renewcommand{\Titulo}{Preparación y aplicación de solución desinfectante para unidades de transporte}

\section{\Titulo}
% \section{Preparación y aplicación de solución desinfectante para unidades de transporte}

\subsection{Objetivo}

\begin{itemize}
	\item Establecer un procedimiento para realizar la preparación de la solución para el desinfectado de las unidades de transporte.
\end{itemize}

\subsection{Alcance}

\begin{itemize}
	\item Solución desinfectante.
\end{itemize}

\subsection{Terminología y definiciones}

\begin{description}
	\defglo{limpieza}
	\defglo{sanitizacion}
\end{description}


\subsection{Procedimiento}

\subsubsection{Materiales}

\begin{itemize}
	\item Preparación de \emph{WHISPER V} a la concentración indicada por el proveedor \qty{0.2}{\percent} o bien (\qty{2}{\milli\liter\per\liter}).
	\item Bote aspersor.
	\item Probeta dosificadora.
\end{itemize}

\subsubsection{Instrucciones}
\paragraph{Desinfección}
\begin{enumerate}
	\item Se procede al preparado de la solución desinfectante.
	\item Agregar en \qty{7}{\liter} de agua dentro del aspersor y posteriormente se adiciona Whisper V \qty{14}{\milli\liter} para alcanzar una concentración de \qty{150}{\milli\gram\per\liter}. Con tiempo de contacto de \qty{60}{\second}.
	\item Se verifica que la unidad venga limpia y libre de cualquier sólido.
	\item Se aplica el producto desinfectante directamente al interior de la unidad (caja) antes de embarcar producto, aplicándolo de la parte más alta a la parte más baja de la caja y de adentro hacia afuera de la unidad.
	\item Se llevará a cabo una inspección diaria. Registrando el indicador y medida correspondiente
	\item Si la desviación se repite frecuentemente se dará curso de capacitación al personal de almacén para que realice eficientemente su trabajo.
	\item Si después de haber capacitado al personal de almacén se siguen presentando desviaciones por causas injustificadas, se tomaran acciones más enérgicas con el personal por incumplimiento con sus deberes.
\end{enumerate}

\subsection{Responsables de la actividad}
\begin{itemize}
	\item \textbf{Ejecutado} por personal de limpieza;
	\item \textbf{Monitoreado} por personal de mantenimiento;
	\item \textbf{Verificado} por personal de Calidad.
\end{itemize}

\subsection{Acciones correctivas}
\begin{itemize}
	\item La tarea es liberada una vez que el área se encuentran libre de polvo, suciedad, residuos de materiales extraños, residuos de jabón, y seca. En caso contrario la tarea se deberá volver a realizar como se indica en el procedimiento.
	\item En caso de no conformidad reportar en \RAC.
\end{itemize}

\subsection{Frecuencia}
\begin{itemize}
	\item Diario de lunes a sábado.
\end{itemize}

\begin{changelog}[simple, sectioncmd=\subsection*,label=changelog-3.8]
	\begin{version}[v=2.0, date=2023--01, author=Pablo E. Alanis]
		\item Cambio de formato;
		\item Cambios en la serialización de versiones;
		\item Correcciones ortográficas y de estilo.
	\end{version}

	\begin{version}[v=1.7, date=2022--05, author=Alonso M.]
		\item cambio de fecha;
	\end{version}

	\shortversion{v=1.6, date=2021--05, changes=No hubo cambios}
\end{changelog}
\renewcommand{\MenorVer}{0}
\renewcommand{\MayorVer}{2}
%\renewcommand{\Edit}{\MayorVer.\MenorVer}
\renewcommand{\Codigo}{HYS-9-IT}
\renewcommand{\FechaPub}{2023--01}
\renewcommand{\Titulo}{Preparación de solución desinfectante}

\section{\Titulo}\index{Preparación de soluciones!desinfectante}

\subsection{Objetivo}
\begin{itemize}
	\item Establecer un procedimiento para realizar la preparación de la solución desinfectante del tapete sanitario.
\end{itemize}

\subsection{Alcance}
\begin{itemize}
	\item Tapete sanitario de aduana de ingreso.
\end{itemize}

\subsection{Terminología y definiciones}

\begin{description}
\defglo{limpieza}
\defglo{sanitizacion}
\end{description}

\subsection{Procedimiento}

\subsubsection{Materiales}
\begin{itemize}
	\item Preparación de \emph{WHISPER V} a la concentración indicada por el proveedor \qty{1.00}{\percent} o bien \qty{10}{\milli\liter\per\liter}.
	\item Cubeta de \qty{20}{\litre}.
	\item Tapete sanitario.
\end{itemize}

\subsubsection{Instrucciones}

\paragraph{Lavado y Desinfección}

\begin{enumerate}
	\item Después del lavado del tapete sanitario se procede al preparado de la solución desinfectante para tapete sanitario (De aduana Sanitaria).
	\item Agregar en \qty{20}{\litre} de agua dentro del tapete sanitario y posteriormente se adiciona Whisper V (\qty{200}{\milli\liter}) para alcanzar una concentración de \qtyrange{400}{1200}{\milli\gram\per\liter}.
	\item Realizar verificaciones de concentración 1 vez por día para garantizar que la solución preparada se encuentra dentro de los valores establecidos.
\end{enumerate}

\subsection{Acciones preventivas}

\begin{itemize}
	\item La tarea es liberada una vez que el área se encuentre libre de polvo, suciedad, residuos de materiales extraños, residuos de jabón, y seca. En caso contrario la tarea se deberá volver a realizar.
	\item Se llevará a cabo una inspección diaria. Registrando el indicador y medida correspondiente
	\item Si la desviación se repite frecuentemente se dará curso de capacitación al personal de limpieza para que realicen eficientemente su trabajo.
	\item Si después de haber capacitado al personal de limpieza se siguen presentando desviaciones por causas injustificadas, se tomaran acciones más enérgicas con el personal por incumplimiento con sus deberes.
\end{itemize}

\subsection{Acciones correctivas}
\begin{itemize}
	\item Cuando se presente una no conformidad en la realización del procedimiento, detectada por el supervisor o encargado del área se deberá de repetir este mismo hasta que sea corregida la desviación
	\item Registrar en formato de acciones correctivas correspondientes.
\end{itemize}

\subsection{Frecuencia}
\begin{itemize}
	\item Diario de lunes a sábado de 08:00 a 08:15.
\end{itemize}

\begin{changelog}[simple, sectioncmd=\subsection*,label=changelog-3.9]
	\begin{version}[v=2.0, date=2023--01, author=Pablo E. Alanis]
		\item Cambio de formato;
		\item Cambios en la serialización de versiones;
		\item Correcciones ortográficas y de estilo.
	\end{version}

	\begin{version}[v=1.7, date=2022--05, author=Alonso M.]
		\item cambio de fecha;
	\end{version}

	\shortversion{v=1.6, date=2021--05, changes=No hubo cambios}
\end{changelog}
\renewcommand{\MenorVer}{0}
\renewcommand{\MayorVer}{2}
%\renewcommand{\Edit}{\MayorVer.\MenorVer}
\renewcommand{\Codigo}{HYS-10-IT}
\renewcommand{\FechaPub}{2023--01}
\renewcommand{\Titulo}{Lavado de patios de carga y colchonetas}

\section{\Titulo}
% \section{Lavado de patios de carga y colchonetas}

\subsection{Objetivo}

\begin{itemize}
	\item Establecer procedimiento de limpieza que se debe de realizar en los patios de carga y colchonetas con el fin de dar mantenimiento de limpieza y minimizar los riesgos de contaminación para el producto.
\end{itemize}

\subsection{Alcance}

\begin{itemize}
	\item Patios de carga y colchonetas.
\end{itemize}

\subsection{Terminología y definiciones}
\begin{description}
\defglo{limpieza}
\defglo{sanitizacion}
\end{description}

%\subsection{Documentos y/o normas relacionadas}

%N/A

\subsection{Procedimiento}

\subsubsection{Materiales}

\begin{itemize}
	\item Preparación de Liquid K a la concentración indicada por el proveedor \qtyrange{0.13}{0.26}{\percent} o bien \qtyrange{1.3}{2.6}{\milli\liter\per\liter}
	\item Cubeta de \qty{20}{\litre}
	\item Cepillo para piso color amarillo
	\item Recogedor de plástico color amarillo
	\item Hidrolavadora
\end{itemize}

\subsubsection{Instrucciones}

\paragraph{Preoperativo}

\subparagraph{Lavado de patios de carga}

\begin{enumerate}
	\item Barrer toda el área con el cepillo para piso.
	\item Deshacer cualquier tipo de desechos en bolsas plásticas o contenedor.
\end{enumerate}

\paragraph{Operativo}

\subparagraph{Limpieza diaria de pisos de carga y exteriores}

\begin{enumerate}
	\item Barrer toda el área con el cepillo para piso.
	\item Desechar cualquier tipo de desechos en bolsas plásticas o contenedor.
	\item Con ayuda de la hidrolavadora eliminar suciedad incrustada en pisos.
	\item En caso de un derrame o suciedad incrustada en el área de patio aplicar desengrasante a la concentración indicada por el proveedor.
	\item Tallar con ayuda del cepillo para pisos.
	\item Retire excedente de detergente con ayuda de la hidrolavadora.
	\item Para finalizar barrer todo el exceso de agua con el cepillo
	\item Elimine excedente de agua
\end{enumerate}

\subparagraph{Lavado de Colchonetas}

\begin{itemize}
	\item Enjuagar con la hidrolavadora asegurándose de eliminar excedentes de polvo,
	\item Posteriormente continuar cepillando el área con detergente líquido (a la concentración indicada por el proveedor).
	\item Con la ayuda de una hidrolavadora con agua a presión retirar el exceso de detergente
	\item Si es requerido seque el área con toalla desechable.
\end{itemize}

\subsection{Responsables de la actividad}

\begin{itemize}
	\item \textbf{Ejecutado} por personal de limpieza
	\item \textbf{Monitoreado} por personal de mantenimiento
	\item \textbf{Verificado} por personal de Calidad.
\end{itemize}

\subsection{Acciones preventivas}

\begin{itemize}
	\item Se llevará a cabo una inspección diaria. Registrando el indicador y medida correspondiente
	\item Si la desviación se repite frecuentemente se dará curso de capacitación al personal de limpieza para que realicen eficientemente su trabajo.
	\item Si después de haber capacitado al personal de limpieza se siguen presentando desviaciones por causas injustificadas, se tomaran acciones más enérgicas con el personal por incumplimiento con sus deberes.
\end{itemize}

\subsection{Acciones correctivas}

\begin{itemize}
	\item La tarea es liberada una vez que el área se encuentran libre de polvo, suciedad, residuos de materiales extraños, residuos de jabón, y seca. En caso contrario la tarea se deberá volver a realizar como se indica en el procedimiento.
	\item En caso de no conformidad reportar en \RAC.
\end{itemize}

\subsection{Frecuencia}

\begin{itemize}
	\item[Superficial] Diario de lunes a viernes de 10:30 a 11:30 y sábados de 11:30 a 12:30;
	\item[Profundo] Trimestralmente.
\end{itemize}

\begin{changelog}[simple, sectioncmd=\subsection*,label=changelog-3.9]
	\begin{version}[v=2.0, date=2023--01, author=Pablo E. Alanis]
		\item Cambio de formato;
		\item Cambios en la serialización de versiones;
		\item Correcciones ortográficas y de estilo.
	\end{version}

	\begin{version}[v=1.7, date=2022--05, author=Alonso M.]
		\item cambio de fecha;
	\end{version}

	\shortversion{v=1.6, date=2021--05, changes=No hubo cambios}
\end{changelog}
\renewcommand{\MenorVer}{0}
\renewcommand{\MayorVer}{2}
%\renewcommand{\Edit}{\MayorVer.\MenorVer}
\renewcommand{\Codigo}{HYS-11-IT}
\renewcommand{\FechaPub}{2023--01}
\renewcommand{\Titulo}{Lavado de patio de montacargas}

\section{\Titulo}\index{Lavado!patio de montacargas, de}
% \section{Lavado de área de montacargas}

\subsection{Objetivo}
Establecer procedimiento de limpieza del área de montacargas.

\subsection{Alcance}
Área de carga de montacargas

\subsection{Terminología y definiciones}
\begin{itemize}
	\item \textbf{Limpieza} Proceso de remover los residuos Orgánicos y suciedad que pueden ser una fuente de contaminación. Los métodos de limpieza adecuados y los materiales dependerán de la naturaleza del alimento.
	\item \textbf{Sanitización} Proceso realizado después de efectuar la limpieza profunda de equipos, maquinaria e instalaciones la cual nos permite el reducir el número de microorganismos.
\end{itemize}

%\subsection{Documentos y/o normas relacionadas}

%N/A

\subsection{Procedimiento}

\subsubsection{Materiales}

\begin{itemize}
	\item Preparación de Liquid K a la concentración indicada por el proveedor (0.13\% a 0.26\% v/v)(1.3ml a 2.6ml en 1L de agua)
	\item Cubeta de \qty{20}{\litre}
	\item Cepillo para piso color rojo
	\item Cepillos manuales color rojo
	\item Recogedor de plástico
	\item Jalador color rojo
	\item Toalla desechable
\end{itemize}

\subsubsection{Instrucciones}

\paragraph{Preoperativo}

\subparagraph{Limpieza de paredes y techos}

\begin{enumerate}
	\item Preparar el área que se va a lavar (retirando la herramienta que se encuentre en el área)
	\item Prepare detergente a la concentración indicada por el proveedor
	\item Aplicar el detergente
	\item Tallar con el cepillo empezando por techos, paredes, dejar reposar 2 a 5 min sin permitir que la espuma se seque
	\item Retire el excedente de la solución iniciando de la parte más alta finalizando en la parte más baja.
	\item Aplicar detergente en el piso
	\item Tallar con cepillo
	\item Retire exceso de humedad con el jalador para piso
\end{enumerate}

\subsection{Responsables de la actividad}

\begin{enumerate}
	\item Ejecutado por personal de limpieza
	\item Monitoreado por personal de mantenimiento
	\item Verificado por personal de Calidad.
\end{enumerate}

\subsection{Acciones preventivas}

\begin{itemize}
	\item Se llevara a cabo una inspección diaria. Registrando el indicador y medida correspondiente
	\item Si la desviación se repite frecuentemente se dará curso de capacitación al personal de limpieza para que realicen eficientemente su trabajo.
	\item Si después de haber capacitado al personal de limpieza se siguen presentando desviaciones por causas injustificadas, se tomaran acciones más enérgicas con el personal por incumplimiento con sus deberes.
\end{itemize}

\subsection{Acciones correctivas}

\begin{itemize}
	\item La tarea es liberada una vez que el área se encuentran libre de polvo, suciedad, residuos de materiales extraños, residuos de jabón, y seca. En caso contrario la tarea se deberá volver a realizar como se indica en el procedimiento.
	\item En caso de no conformidad reportar en \RAC.
\end{itemize}

\subsection{Frecuencia}

\begin{itemize}
	\item Diario de lunes a viernes de 13:30pm a 14:30pm y sábados de 10:30am a 11:30am / Trimestral revisión del programa general de limpieza.
\end{itemize}

\subsection{Historial de modificaciones}

\begin{itemize}
	\item Cuarta edición: cambio de fecha de 24 de mayo 2017 a 28 de enero 2019, se realizó cambio de código de HS-P-LAM a PR-029 Y no se realizaron cambios de la revisión 003 a la 004.
	\item Quinta edición: febrero 2020 se hizo cambio de formato y cambio de código de PR-029 a IT-OP-012. Y no se realizaron cambios de la revisión 04 a la 05.
	\item Sexta edición: febrero 2021 no se realizaron cambios de la revisión 05 a la 06.
	\item Séptima edición: febrero 2022 no se realizaron cambios de la revisión 06 a la 07.
\end{itemize}

\thispagestyle{formato-PI}
\renewcommand{\MenorVer}{0}
\renewcommand{\MayorVer}{2}
%\renewcommand{\Edit}{\MayorVer.\MenorVer}
\renewcommand{\Codigo}{HYS-12-IT}
\renewcommand{\FechaPub}{2023--01}
\renewcommand{\Titulo}{Limpieza de accesorios de limpieza}

\section{\Titulo}\index{Limpieza!accesorios de limpieza, de}
% \section{Limpieza de accesorios de limpieza}

\subsection{Objetivo}
Establecer procedimiento que aseguren un mantenimiento y limpieza apropiada de utensilios de limpieza.

\subsection{Alcance}
Utensilios de limpieza del área de almacén.

\subsection{Terminología y definiciones}

\begin{description}
	\defglo{limpieza}
	\defglo{sanitizacion}
\end{description}

\subsection{Procedimiento}
\subsubsection{Materiales}
\begin{itemize}
	\item Preparación de Liquid K a la concentración indicada por el proveedor,\qtyrange{0.13}{0.26}{\percent}, o bien \qtyrange{1.3}{2.6}{\milli\liter\per\liter}.
	\item Cubeta de \qty{20}{\litre}
\end{itemize}

\subsubsection{Instrucciones}

\paragraph{Operativo}

\subparagraph{Limpieza de accesorios de limpieza}

\begin{enumerate}
	\item Enjuagar los utensilios como, cepillo de mano, recogedores, cubetas, jaladores que se utilizaron en el transcurso del día.
	\item Sumergir en una cubeta con agua y detergente
	\item Enjuagarlos perfectamente
	\item Elimine el excedente de humedad
	\item Colocarlas en su lugar correspondiente para terminar de secar.
	\item Lavar cepillos, jaladores y recogedores de la parte de arriba hacia abajo y enjuagar perfectamente.
	\item Después enjuagar y colocarlos en el lugar correspondiente.
	\item Enjuague perfectamente cubetas hasta eliminar residuos de suciedad.
	\item De igual manera, elimine el excedente de humedad.
	\item[\textbf{Nota 1:}] NO permita la acumulación de agua en cubetas o utensilios de limpieza
	\item[\textbf{Nota 2:}] La limpieza de los utensilios se realiza de acuerdo a su mismo código de color.
\end{enumerate}

\subsection{Responsables de la actividad}
\begin{itemize}
	\item \textbf{Ejecutado} por personal de limpieza;
	\item \textbf{Monitoreado} por personal de mantenimiento;
	\item \textbf{Verificado} por personal de Calidad.
\end{itemize}

\subsection{Acciones preventivas}
\begin{itemize}
	\item Se llevará a cabo una inspección diaria, registrando el indicador y medida correspondiente
	\item Si la desviación se repite frecuentemente se dará curso de capacitación al personal de limpieza para que realicen eficientemente su trabajo.
	\item Si después de haber capacitado al personal de limpieza se siguen presentando desviaciones por causas injustificadas, se tomaran acciones más enérgicas con el personal por incumplimiento con sus deberes.
\end{itemize}

\subsection{Acciones correctivas}
\begin{itemize}
	\item La tarea es liberada una vez que el área se encuentran libre de polvo, suciedad, residuos de materiales extraños, residuos de jabón, y seca. En caso contrario la tarea se deberá volver a realizar como se indica en el procedimiento.
	\item En caso de no conformidad reportar en \RAC.
\end{itemize}

\subsection{Frecuencia}
\begin{itemize}
	\item Diario de lunes a viernes de 16:00 a 16:30 y sábados de 11:30 a 12:00.
\end{itemize}

\begin{changelog}[title=Registro de cambios,simple, sectioncmd=\subsection*,label=changelog-\thesection-\MayorVer.\MenorVer2]
	\begin{version}[v=\MayorVer.\MenorVer, date=2023--01, author=Pablo E. Alanis]
		\item Cambio de formato;
		\item Cambios en la serialización de versiones;
		\item Correcciones ortográficas y de estilo.
	\end{version}

	\begin{version}[v=1.7, date=2022--05, author=Alonso M.]
		\item cambio de fecha;
	\end{version}

	\shortversion{v=1.6, date=2021--05, changes=No hubo cambios}
\end{changelog}
\thispagestyle{formato-PI}
\renewcommand{\MenorVer}{0}
\renewcommand{\MayorVer}{2}
%\renewcommand{\Edit}{\MayorVer.\MenorVer}
\renewcommand{\Codigo}{HYS-13-IT}
\renewcommand{\FechaPub}{2023--01}
\renewcommand{\Titulo}{Limpieza de utensilios de limpieza de alérgenos}

\section{\Titulo}\index{Limpieza de utensilios!limpieza de alérgenos, de}
% \section{Limpieza de accesorios de alérgenos}

\subsection{Objetivo}
Establecer un procedimiento que aseguren un mantenimiento y limpieza apropiada en caso de presentarse un derrame de producto Alérgeno.

\subsection{Alcance}
Toda superficie afectada por un derrame de alimento con componente(s) alérgenico(s).

\subsection{Terminología y definiciones}
\begin{itemize}
	\item \textbf{Limpieza:} Proceso de remover los residuos de Orgánicos y suciedad que pueden ser una fuente de contaminación. Los métodos de limpieza adecuados y los materiales dependerán de la naturaleza del alimento.
	\item \textbf{Sanitización:} Proceso realizado después de efectuar la limpieza profunda de equipos, maquinaria e instalaciones la cual nos permite el reducir el número de microorganismos
\end{itemize}
\subsection{Procedimiento}
\subsubsection{Materiales}
\begin{itemize}
	\item Preparación de \emph{Liquid K} a la concentración indicada por el proveedor, \qtyrange{0.13}{0.26}{\percent} o bien \qtyrange{1.3}{2.6}{\milli\liter\per\liter};
	\item Cubeta de \qty{20}{\litre};
	\item Cepillo para piso color Blanco;
	\item Cepillos manuales color Blanco;
	\item Recogedor de plástico;
	\item Jalador color Blanco;
	\item Toallas desechables.
\end{itemize}

\subsubsection{Instrucciones}
\paragraph{Derrame de alérgenos en el piso de almacén, cámaras frías y andenes}

\begin{enumerate}
	\item Verificar la no afectación de ningún otro producto
	\item Recolectar el producto alérgeno de las áreas afectadas con ayuda de toallas desechables (elimine toallas después de su uso)
	\item Prepare detergente a la concentración indicada por el proveedor.
	\item Tallar el piso con agua y jabón y el cepillo asignado para estos casos
	\item Enjuagar el piso
	\item Eliminar la humedad con el jalador asignado
	\item Dar aviso al supervisor de Almacén y Calidad para realizar revisión visual
	\item Registro de lavado por derrame de material alérgeno.
	\item Si se afectó algún otro producto o empaque se procede a eliminar el producto contaminante por medio de limpieza en seco (toallas desechables).
	\item Separar el producto y etiquetar como producto detenido y/o retenido
	\item Valorar el grado de afectación del producto en coordinación con el Coordinador de calidad del cliente en cuestión.
	\item Se precederá según indicaciones sobre la disposición del producto.
	\item Registro de lavado por derrame de material alérgeno.
\end{enumerate}

\subsection{Responsables de la actividad}
\begin{itemize}
	\item \textbf{Ejecutado} por personal de limpieza;
	\item \textbf{Monitoreado} por personal de mantenimiento;
	\item \textbf{Verificado} por personal de Calidad.
\end{itemize}

\subsection{Acciones preventivas}
\begin{itemize}
	\item Se llevará a cabo una inspección diaria. Registrando el indicador y medida correspondiente
	\item Si la desviación se repite frecuentemente se dará curso de capacitación al personal de almacén y limpieza para que realicen eficientemente su trabajo.
	\item Si después de haber capacitado al personal de almacén y limpieza se siguen presentando desviaciones por causas injustificadas, se tomaran acciones más enérgicas con el personal por incumplimiento con sus deberes.
\end{itemize}

\subsection{Acciones correctivas}
\begin{enumerate}
	\item La tarea es liberada una vez que el área se encuentran libre de polvo, suciedad, residuos de materiales extraños, residuos de jabón, y seca. En caso contrario la tarea se deberá volver a realizar como se indica en el procedimiento.
	\item En caso de no conformidad reportar en \RAC.
\end{enumerate}

\subsection{Frecuencia}
Cuando se presente la incidencia.

\begin{changelog}[simple, sectioncmd=\subsection*,label=changelog-3.13]
	\begin{version}[v=2.0, date=2023--01, author=Pablo E. Alanis]
		\item Cambio de formato;
		\item Cambios en la serialización de versiones;
		\item Correcciones ortográficas y de estilo.
	\end{version}

	\begin{version}[v=1.7, date=2022--05, author=Alonso M.]
		\item cambio de fecha;
	\end{version}	

	\shortversion{v=1.6, date=2021--05, changes=No hubo cambios}
\end{changelog}
\thispagestyle{formato-PI}
\renewcommand{\MenorVer}{0}
\renewcommand{\MayorVer}{2}
%\renewcommand{\Edit}{\MayorVer.\MenorVer}
\renewcommand{\Codigo}{HYS-14-IT}
\renewcommand{\FechaPub}{2023--01}
\renewcommand{\Titulo}{Limpieza de oficinas}

\section{\Titulo}\index{Limpieza!oficinas, de}
% \section{Limpieza de oficinas}

\subsection{Objetivo}
Establecer un procedimiento que aseguren un mantenimiento y limpieza apropiada de Oficinas.

\subsection{Alcance}
Todas las oficinas del almacén.

\subsection{Terminología y definiciones}
\begin{description}
\defglo{limpieza}
\defglo{sanitizacion}
\end{description}

\subsection{Procedimiento}
\subsubsection{Materiales}
\begin{itemize}
	\item Preparación de \emph{Liquid K} a la concentración indicada por el proveedor \qtyrange{.13}{.26}{\percent}, o bien \qtyrange{1.3}{2.6}{\milli\liter\per\liter}.
	\item Cubeta de \qty{20}{\litre}
	\item Cepillo para piso color verde
	\item Recogedor de plástico color verde
	\item Toalla desechable
	\item Trapeador color verde
\end{itemize}

\subsubsection{Instrucciones}
\paragraph{Limpieza de pisos}
\begin{enumerate}
	\item Barrer pisos de áreas de oficinas.
	\item Recoger la suciedad recolectada con recogedor.
	\item Colocar residuos en bolsa para basura.
	\item Prepara cubeta con agua y detergente a la concentración indicada por el proveedor para iniciar el trapeado de áreas.
	\item Trapear las áreas de oficinas.
	\item Enjuagar trapeador con agua limpia.
	\item Eliminar excedente de humedad.
\end{enumerate}

\paragraph{Limpieza de escritorios}

\begin{enumerate}
	\item Preparar disolución a dosificación indicada
	\item Humedecer una toalla desechable
	\item Realizar limpieza de superficie y base de escritorio
	\item Eliminar excedente de humedad
\end{enumerate}

\paragraph{Limpieza de ventanas:}
\begin{enumerate}
	\item Preparar disolución a dosificación indicada por el proveedor
	\item Humedecer una toalla desechable
	\item Realizar limpieza de superficie
	\item Eliminar excedente de humedad
\end{enumerate}

\subsection{Responsables de la actividad}
\begin{itemize}
	\item \textbf{Ejecutado} por personal de limpieza;
	\item \textbf{Monitoreado} por personal de mantenimiento;
	\item \textbf{Verificado} por personal de Calidad.
\end{itemize}

\subsection{Acciones preventivas}
\begin{itemize}
	\item Se llevará a cabo una inspección diaria. Registrando el indicador y medida correspondiente
	\item Si la desviación se repite frecuentemente se dará curso de capacitación al personal de limpieza para que realicen eficientemente su trabajo.
	\item Si después de haber capacitado al personal de limpieza se siguen presentando desviaciones por causas injustificadas, se tomaran acciones más enérgicas con el personal por incumplimiento con sus deberes.
\end{itemize}

\subsection{Acciones correctivas}

\begin{itemize}
	\item La tarea es liberada una vez que el área se encuentran libre de polvo, suciedad, residuos de materiales extraños, residuos de jabón, y seca. En caso contrario la tarea se deberá volver a realizar como se indica en el procedimiento.
	\item En caso de no conformidad reportar en \RAC.
\end{itemize}

\subsection{Frecuencia}
Diario de lunes a sábado de 8:00 a 10:30.

\begin{changelog}[title=Registro de cambios,simple, sectioncmd=\subsection*,label=changelog-\thesection-\MayorVer.\MenorVer4]
	\begin{version}[v=\MayorVer.\MenorVer, date=2023--01, author=Pablo E. Alanis]
		\item Cambio de formato;
		\item Cambios en la serialización de versiones;
		\item Correcciones ortográficas y de estilo.
	\end{version}

	\begin{version}[v=1.7, date=2022--05, author=Alonso M.]
		\item cambio de fecha;
	\end{version}	

	\shortversion{v=1.6, date=2021--05, changes=No hubo cambios}
\end{changelog}
\thispagestyle{formato-PI}
\renewcommand{\MenorVer}{0}
\renewcommand{\MayorVer}{2}
%\renewcommand{\Edit}{\MayorVer.\MenorVer}
\renewcommand{\Codigo}{HYS-15-IT}
\renewcommand{\FechaPub}{2023--01}
\renewcommand{\Titulo}{Limpieza de contenedores de basura}

\section{\Titulo}
% \section{Limpieza de contenedores de basura}

\subsection{Objetivo}
Establecer un procedimiento que aseguren un mantenimiento y limpieza apropiada de contenedor de basura.

\subsection{Alcance}
Contenedores de basura dentro de almacén.

\subsection{Terminología y definiciones}
\begin{description}
\defglo{limpieza}
\defglo{sanitizacion}
\end{description}

\subsection{Procedimiento}

\subsubsection{Materiales}
\begin{itemize}
	\item Preparación de Liquid K a la concentración indicada por el proveedor, \qtyrange{.13}{.26}{\percent}, o bien \qtyrange{1.3}{2.6}{\milli\liter\per\liter}
	\item Cubeta de \qty{20}{\litre}
	\item Cepillo color amarillo
	\item Guantes para agua
\end{itemize}

\subsubsection{Instrucciones}

\begin{enumerate}
	\item Retirar la basura del bote o contenedor.
	\item Enjuagar con agua.
	\item Aplicar el detergente a la concentración indicada por el proveedor,
	\item Con el cepillo para piso tallar la pared del cesto de arriba hacia abajo, por la parte de adentro y tallar la parte de abajo por dentro.
	\item Tallar las paredes externas, de arriba hacia abajo, después inclinarlo y tallar la base externa.
	\item Enjuagar por la parte de afuera y después la parte de adentro.
	\item Hacer una inspección visual.
\end{enumerate}

\subsection{Responsables de la actividad}

\begin{itemize}
	\item \textbf{Ejecutado} por personal de limpieza;
	\item \textbf{Monitoreado} por personal de mantenimiento;
	\item \textbf{Verificado} por personal de Calidad.
\end{itemize}

\subsection{Acciones preventivas}

\begin{itemize}
	\item Se llevará a cabo una inspección diaria. Registrando el indicador y medida correspondiente
	\item Si la desviación se repite frecuentemente se dará curso de capacitación al personal de limpieza para que realicen eficientemente su trabajo.
	\item Si después de haber capacitado al personal de limpieza se siguen presentando desviaciones por causas injustificadas, se tomaran acciones más enérgicas con el personal por incumplimiento con sus deberes.
\end{itemize}

\subsection{Acciones correctivas}

\begin{itemize}
	\item La tarea es liberada una vez que el área se encuentran libre de polvo, suciedad, residuos de materiales extraños, residuos de jabón, y seca. En caso contrario la tarea se deberá volver a realizar como se indica en el procedimiento.
	\item En caso de no conformidad reportar en \RAC.
\end{itemize}

\subsection{Frecuencia}

Mensual revisión del programa general de limpieza.

\begin{changelog}[simple, sectioncmd=\subsection*,label=changelog-\thesection-\MayorVer.\MenorVer5]
	\begin{version}[v=\MayorVer.\MenorVer, date=2023--01, author=Pablo E. Alanis]
		\item Cambio de formato;
		\item Cambios en la serialización de versiones;
		\item Correcciones ortográficas y de estilo.
	\end{version}

	\begin{version}[v=1.7, date=2022--05, author=Alonso M.]
		\item cambio de fecha;
	\end{version}	

	\shortversion{v=1.6, date=2021--05, changes=No hubo cambios}
\end{changelog}
\thispagestyle{formato-PI}
\renewcommand{\MenorVer}{0}
\renewcommand{\MayorVer}{2}
%\renewcommand{\Edit}{\MayorVer.\MenorVer}
\renewcommand{\Codigo}{HYS-16-IT}
\renewcommand{\FechaPub}{2023--01}
\renewcommand{\Titulo}{Limpieza de baños de empleados}

\section{\Titulo}\index{Lavado!baño de empleados, de}
% \section{Limpieza de baños de empleados}

\subsection{Objetivo}
Establecer un procedimiento que aseguren un mantenimiento y limpieza apropiada de los baños para colaboradores de la empresa.

\subsection{Alcance}

Todos los baños dentro de las instalaciones.

\subsection{Terminología y definiciones}

\begin{description}
\defglo{limpieza}
\defglo{sanitizacion}
\end{description}

%\subsection{Documentos y/o normas relacionadas}

%N/A

\subsection{Procedimiento}

\subsubsection{Materiales}

\begin{itemize}
	\item Preparación de Liquid K a la concentración indicada por el proveedor, \qtyrange{.13}{.26}{\percent}, o bien \qtyrange{1.3}{2.6}{\milli\liter\per\liter};
	\item Cubeta de \qty{20}{\litre};
	\item Cepillo para piso color azul;
	\item Cepillos manuales color azul;
	\item Recogedor de plástico;
	\item Jalador color azul;
	\item Trapeador color azul.
\end{itemize}

\subsubsection{Instrucciones}
\paragraph{Operativo}
\subparagraph{Mingitorios, Sanitarios, lavamanos y pisos}
\begin{enumerate}
	\item Preparar detergente a la concentración indicada por el proveedor
	\item Aplicar detergente en lavamanos
	\item Tallar con el cepillo de mano.
	\item Aplicar detergente a Mingitorios y sanitarios
	\item Tallar con el cepillo identificado para esta área.
	\item Enjuagar con agua limpia.
	\item Eliminar el excedente de agua hacia una coladera.
	\item Aplique detergente en pisos y talle con el cepillo para piso.
	\item Elimine excedente de agua.
\end{enumerate}

\subsection{Responsables de la actividad}

\begin{itemize}
	\item \textbf{Ejecutado} por personal de limpieza;
	\item \textbf{Monitoreado} por personal de mantenimiento;
	\item \textbf{Verificado} por personal de Calidad.
\end{itemize}

\subsection{Acciones preventivas}

\begin{itemize}
	\item Se llevará a cabo una inspección diaria. Registrando el indicador y medida correspondiente
	\item Si la desviación se repite frecuentemente se dará curso de capacitación al personal de limpieza para que realicen eficientemente su trabajo
	\item Si después de haber capacitado al personal de limpieza se siguen presentando desviaciones por causas injustificadas, se tomaran acciones más enérgicas con el personal por incumplimiento con sus deberes.
\end{itemize}

\subsection{Acciones correctivas}

\begin{itemize}
	\item La tarea es liberada una vez que el área se encuentran libre de polvo, suciedad, residuos de materiales extraños, residuos de jabón, y seca. En caso contrario la tarea se deberá volver a realizar como se indica en el procedimiento.
	\item En caso de no conformidad reportar en \RAC.
\end{itemize}

\subsection{Frecuencia}

Diario de lunes a viernes de 10:30 a 11:00 y sábados de 11:30 a 12:00.

\begin{changelog}[simple, sectioncmd=\subsection*,label=changelog-3.14]
	\begin{version}[v=2.0, date=2023--01, author=Pablo E. Alanis]
		\item Cambio de formato;
		\item Cambios en la serialización de versiones;
		\item Correcciones ortográficas y de estilo.
	\end{version}

	\begin{version}[v=1.7, date=2022--05, author=Alonso M.]
		\item cambio de fecha;
	\end{version}	

	\shortversion{v=1.6, date=2021--05, changes=No hubo cambios}
\end{changelog}
\thispagestyle{formato-PI}
\renewcommand{\MenorVer}{0}
\renewcommand{\MayorVer}{2}
%\renewcommand{\Edit}{\MayorVer.\MenorVer}
\renewcommand{\Codigo}{HYS-17-IT}
\renewcommand{\FechaPub}{2023--01}
\renewcommand{\Titulo}{Lavado del comedor}

\section{\Titulo}\index{Limpieza!comedor, del}
% \section{Lavado del comedor}

\subsection{Objetivo}
Establecer un procedimiento que aseguren un mantenimiento y limpieza apropiada del comedor.

\subsection{Alcance}
Comedor de almacén.

\subsection{Terminología y definiciones}
\begin{description}
\defglo{limpieza}
\defglo{sanitizacion}
\end{description}

%\subsection{Documentos y/o normas relacionadas}

%N/A

\subsection{Procedimiento}

\subsubsection{Materiales}

\begin{itemize}
	\item Preparación de Liquid K a la concentración indicada por el proveedor \qtyrange{.13}{.26}{\percent}, o bien \qtyrange{1.3}{2.6}{\milli\liter\per\liter}
	\item Cubeta de \qty{20}{\litre}
	\item Cepillo para piso color verde
	\item Fibra
	\item Toalla desechable
	\item Recogedor de plástico color verde
	\item Jalador color verde
	\item Trapeador color verde
\end{itemize}

\subsubsection{Instrucciones}
\paragraph{Operativo}
\subparagraph{Limpieza para mesa y equipo de comedor}
\begin{enumerate}
	\item Eliminar toda materia orgánica e inorgánica que se encuentre en estas aéreas y depositarla en los contenedores de basura
	\item Desconectar todos los equipos eléctricos.
	\item Preparar detergente a la concentración indicada por el proveedor
	\item Aplicar detergente, tallar con fibra la mesa y el equipo de comedor
	\item Enjuagar con agua limpia
	\item Eliminar excedente de humedad con una toalla desechable
\end{enumerate}

\subparagraph{Limpieza de pisos}
\begin{enumerate}
	\item Eliminar toda materia orgánica e inorgánica que se encuentre en estas aéreas y depositarlas en los contenedores de basura
	\item Prepare detergente a la concentración indicada por el proveedor
	\item Con ayuda de trapeador realizar limpieza del piso
	\item Enjuague el trapeador con agua limpia.
	\item Con el trapeador elimine el excedente de humedad.
\end{enumerate}

\subsection{Responsables de la actividad}

\begin{itemize}
	\item \textbf{Ejecutado} por personal de limpieza;
	\item \textbf{Monitoreado} por personal de mantenimiento;
	\item \textbf{Verificado} por personal de Calidad.
\end{itemize}

\subsection{Acciones preventivas}

\begin{itemize}
	\item Se llevará a cabo una inspección diaria. Registrando el indicador y medida correspondiente
	\item Si la desviación se repite frecuentemente se dará curso de capacitación al personal de limpieza para que realicen eficientemente su trabajo.
	\item Si después de haber capacitado al personal de limpieza se siguen presentando desviaciones por causas injustificadas, se tomaran acciones más enérgicas con el personal por incumplimiento con sus deberes.
\end{itemize}

\subsection{Acciones correctivas}

\begin{itemize}
	\item La tarea es liberada una vez que el área se encuentran libre de polvo, suciedad, residuos de materiales extraños, residuos de jabón, y seca. En caso contrario la tarea se deberá volver a realizar como se indica en el procedimiento.
	\item En caso de no conformidad reportar en \RAC.
\end{itemize}

\subsection{Frecuencia}

Diaria de lunes a viernes de 14:30 a 15:00 y sábados de 11:30 a 12:00.

\begin{changelog}[title=Registro de cambios,simple, sectioncmd=\subsection*,label=changelog-\thesection-\MayorVer.\MenorVer7]
	\begin{version}[v=\MayorVer.\MenorVer, date=2023--01, author=Pablo E. Alanis]
		\item Cambio de formato;
		\item Cambios en la serialización de versiones;
		\item Correcciones ortográficas y de estilo.
	\end{version}

	\begin{version}[v=1.7, date=2022--05, author=Alonso M.]
		\item cambio de fecha;
	\end{version}	

	\shortversion{v=1.6, date=2021--05, changes=No hubo cambios}
\end{changelog}
\thispagestyle{formato-PI}
\renewcommand{\MenorVer}{0}
\renewcommand{\MayorVer}{2}
%\renewcommand{\Edit}{\MayorVer.\MenorVer}
\renewcommand{\Codigo}{HYS-18-IT}
\renewcommand{\FechaPub}{2023--01}
\renewcommand{\Titulo}{Limpieza de uniformes}

\section{\Titulo}\index{Limpieza!uniformes, de}
% \section{Limpieza de uniformes}

\subsection{Objetivo}
Asegurar el lavado de los uniformes que se utilizan en el almacén.

\subsection{Alcance}
Uniformes de trabajo.

\subsection{Terminología y definiciones}
\begin{description}
\defglo{limpieza}
\defglo{sanitizacion}
\end{description}

\subsection{Procedimiento}
\subsubsection{Materiales}
\begin{itemize}
	\item Bitácora de registros.
\end{itemize}

\subsubsection{Instrucciones}
\paragraph{Limpieza de uniformes}
\begin{enumerate}
	\item El personal dejara los uniformes con el personal de recepción o vigilancia firmara de entrega de uniforme, cantidad y la fecha.
	\item El personal entregara los uniformes al personal de lavandería el cual pasara a las instalaciones.
	\item La Empresa Externa (Servicio de Lavado) es responsable de lo siguiente:
	\item Brindar el servicio de lavado de buena calidad de la empresa \gls{RDF}.
	\item Verificar condiciones y cantidades de uniformes que la empresa le hace entrega, así como empleado, departamento y firma de entregado y recibido.
	\item Entregar el uniforme lavado en las fechas pactadas.
	\item Entregar las notas de remisión correspondientes a cada entrega realizada.
	\item Se entregara uniformes al personal y este firmara de acuerdo de recibido.
\end{enumerate}

\subsection{Responsables de la actividad}

\begin{itemize}
	\item \textbf{Ejecutado} por personal de limpieza;
	\item \textbf{Monitoreado} por personal de mantenimiento;
	\item \textbf{Verificado} por personal de Calidad.
\end{itemize}

\subsection{Acciones preventivas}

\begin{itemize}
	\item Se llevará a cabo una inspección diaria. Registrando el indicador y medida correspondiente
	\item Si la desviación se repite frecuentemente se dará curso de capacitación al personal de limpieza para que realicen eficientemente su trabajo.
	\item Si después de haber capacitado al personal de limpieza se siguen presentando desviaciones por causas injustificadas, se tomaran acciones más enérgicas con el personal por incumplimiento con sus deberes.
\end{itemize}

\subsection{Acciones correctivas}

\begin{itemize}
	\item Cuando se presente una desviación en la realización del procedimiento detectada por el personal se reportara a Gerencia.
\end{itemize}

\subsection{Frecuencia}

\begin{itemize}
	\item Quincenalmente.
\end{itemize}

\begin{changelog}[title=Registro de cambios,simple, sectioncmd=\subsection*,label=changelog-\thesection-\MayorVer.\MenorVer8]
	\begin{version}[v=\MayorVer.\MenorVer, date=2023--01, author=Pablo E. Alanis]
		\item Cambio de formato;
		\item Cambios en la serialización de versiones;
		\item Correcciones ortográficas y de estilo.
	\end{version}

	\begin{version}[v=1.7, date=2022--05, author=Alonso M.]
		\item cambio de fecha;
	\end{version}	

	\shortversion{v=1.6, date=2021--05, changes=No hubo cambios}
\end{changelog}
\thispagestyle{formato-PI}
\renewcommand{\MenorVer}{0}
\renewcommand{\MayorVer}{2}
%\renewcommand{\Edit}{\MayorVer.\MenorVer}
\renewcommand{\Codigo}{HYS-19-IT}
\renewcommand{\FechaPub}{2023--01}
\renewcommand{\Titulo}{Verificación de la eficacia de POES}

\section{\Titulo}\index{Verificación!eficacia de POES, de la}
% \section{Verificación de eficacia de POES}

\subsection{Objetivo}

Establecer los parámetros con los cuales se mide la eficacia de los procedimientos de opresión estándar de sanitización con los cuales se toma la decisión si el procedimiento se realizó de manera adecuada.

\subsection{Alcance}

El alcance es aplicado para todos los procedimientos de operación estándar de sanitización.

\subsection{Terminología y definiciones}

\begin{description}
\defglo{limpieza}
\defglo{sanitizacion}
\end{description}

%\subsection{Documentos y/o normas relacionadas}

%N/A

\subsection{Procedimiento}

\subsubsection{Materiales}

\begin{itemize}
	\item Lámpara
\end{itemize}

\subsubsection{Instrucciones}
\paragraph{Revisión visual de los equipos y áreas posterior a la realización de los procedimientos de operación estándar de sanitización}
\begin{itemize}
	\item El área debe de estar libre de polvo.
	\item El área debe de estar libre de astillas de madera generadas por las tarimas.
	\item El área debe de estar libre de emplaye.
	\item El área debe estar libre de material orgánico o suciedad.
	\item El área debe de estar libre de cualquier material ajeno a la construcción del área o equipo (tornillos, clavos, vidrios, etc.)
\end{itemize}

\subsection{Responsables de la actividad}
\begin{itemize}
	\item \textbf{Verificado} por personal de Calidad.
\end{itemize}

\subsection{Acciones preventivas}
La acción preventiva se lleva a cabo al tiempo de realizar la verificación del cumplimiento de los procedimientos de operación estándar de sanitización.

\subsection{Acciones correctivas}
\begin{itemize}
	\item La tarea es liberada una vez que el área se encuentran libre de polvo, suciedad, residuos de materiales extraños, residuos de jabón, y seca. En caso contrario la tarea se deberá volver a realizar como se indica en el procedimiento.
	\item En caso de no conformidad reportar en \RAC.
\end{itemize}

\subsection{Frecuencia}
En cada verificación de tarea realizada

\begin{changelog}[title=Registro de cambios,simple, sectioncmd=\subsection*,label=changelog-\thesection-\MayorVer.\MenorVer9]
	\begin{version}[v=\MayorVer.\MenorVer, date=2023--01, author=Pablo E. Alanis]
		\item Cambio de formato;
		\item Cambios en la serialización de versiones;
		\item Correcciones ortográficas y de estilo.
	\end{version}

	\begin{version}[v=1.5, date=2022--05, author=Alonso M.]
		\item cambio de fecha;
	\end{version}	

	\shortversion{v=1.4, date=2021--05, changes=No hubo cambios}
\end{changelog}

\thispagestyle{formato-PI}
\renewcommand{\MenorVer}{0}
\renewcommand{\MayorVer}{2}
%\renewcommand{\Edit}{\MayorVer.\MenorVer}
\renewcommand{\Codigo}{HYS-20-IT}
\renewcommand{\FechaPub}{2023--01}
\renewcommand{\Titulo}{Limpieza de manos}

\section{\Titulo}\index{Lavado!manos, de}
% \section{Limpieza de manos}

\subsection{Objetivo}
\begin{itemize}
	\item Establecer el procedimiento del correcto lavado de manos.
\end{itemize}

\subsection{Alcance}
\begin{itemize}
	\item Todo personal y visitante de \gls{RDF}.
\end{itemize}

\subsection{Terminología y definiciones}
\begin{description}
\defglo{limpieza}
\defglo{sanitizacion}
\end{description}

\subsection{Procedimiento}
\subsubsection{Materiales}
\begin{enumerate}
	\item Jabón para manos.
	\item Toallas desechables
	\item desinfectante de manos
\end{enumerate}

\subsubsection{Instrucciones}

\begin{enumerate}
	\item Moje sus manos con agua. En algunos casos deberá accionar el flujo de agua con la con su rodilla o pie.
	\item Coloque suficiente jabón para cubrir toda la superficie de manos y dedos.
	\item Con movimientos circulares frote las palmas de su mano contra la otra.
	\item Con movimientos ascendentes frote el dorso de ambas manos y entre los dedos.
	\item Entrelace sus dedos y con movimientos ascendentes y descendentes talle entre sus dedos.
	\item Entrelace sus dedos y frote para lavar la parte frontal de sus dedos.
	\item Con la ayuda de sus palmas y movimientos giratorios, talle sus dedos pulgares.
	\item Sobre la palma de sus manos, frote la punta de los dedos de la mano contraria.
	\item Enjuague con agua hasta eliminar por completo el jabón de sus manos.
	\item Con ayuda de toallas desechables o secador de aire caliente elimine el agua de sus manos.
	\item Coloque suficiente gel desinfectante para cubrir manos y dedos, frote hasta que este se evapore por completo.
	\item Sus manos se encuentran limpias para poder ingresar o regresar a sus áreas de trabajo.
\end{enumerate}

\subsection{Responsables de la actividad}

\begin{itemize}
	\item \textbf{Ejecutado} por todo el personal y visitantes.
\end{itemize}

\subsection{Acciones preventivas}
\begin{itemize}
	\item La acción preventiva se lleva a cabo al tiempo de realizar la verificación del cumplimiento de los procedimientos de operación estándar de sanitización.
\end{itemize}

\subsection{Acciones correctivas}

\begin{itemize}
	\item La tarea es liberada una vez que el área se encuentran libre de polvo, suciedad, residuos de materiales extraños, residuos de jabón, y seca. En caso contrario la tarea se deberá volver a realizar como se indica en el procedimiento;
	\item En caso de no conformidad reportar en \RAC.
\end{itemize}

\subsection{Frecuencia}

\begin{itemize}
	\item Diaria, antes de ingresar al almacén, en el cambio de actividades, después de ir al baño, antes de comidas.
\end{itemize}

\begin{changelog}[simple, sectioncmd=\subsection*,label=changelog-3.20]
	\begin{version}[v=2.0, date=2023--01, author=Pablo E. Alanis]
		\item Cambio de formato;
		\item Cambios en la serialización de versiones;
		\item Correcciones ortográficas y de estilo.
	\end{version}

	\begin{version}[v=1.5, date=2022--05, author=Alonso M.]
		\item cambio de fecha;
	\end{version}	

	\shortversion{v=1.4, date=2021--05, changes=No hubo cambios}
\end{changelog}
\thispagestyle{formato-PI}
\renewcommand{\MenorVer}{0}
\renewcommand{\MayorVer}{2}
%\renewcommand{\Edit}{\MayorVer.\MenorVer}
\renewcommand{\Codigo}{HYS-21-IT}
\renewcommand{\FechaPub}{2023--01}
\renewcommand{\Titulo}{Determinación de cloruro en agua}

\section{\Titulo}\index{Determinación!cloruro en agua, de}
% \section{Determinación de cloración en agua}

\subsection{Objetivo}
\begin{itemize}
	\item Establecer el procedimiento de como determinar la cantidad de cloro en el agua.
\end{itemize}

\subsection{Alcance}
\begin{itemize}
	\item Este procedimiento de determinación de cloruro es aplicable a todos los puntos de tuberías de tomas de agua ubicadas en las instalaciones de \gls{RDF}.
\end{itemize}

\subsection{Terminología y definiciones}
\begin{description}
\defglo{limpieza}
\defglo{sanitizacion}
\end{description}

\subsection{Procedimiento}
\subsubsection{Precauciones de seguridad}
\begin{itemize}
	\item El indicador empleado para la determinación de cloro, 2-metilanilina, es cancerígeno. Se deben de seguir las indicaciones del MSDS y evitar la exposición directa a este compuesto.
\end{itemize}

\subsubsection{Materiales}
\begin{enumerate}
	\item Indicador (2-metilanilina);
	\item Indicados de concentración de cloro;
	\item Vial de muestra para agua.
\end{enumerate}

\subsubsection{Instrucciones}

\begin{enumerate}
	\item Deje correr agua por un promedio de \qtyrange{1.5}{2}{\minute} para poder tomar una muestra.
	\item Posterior a este tiempo prosiga a llenar el vial del indicador de concentración de cloro.
	\item Vierta 5 gotas del indicador ortotolidina en la muestra de agua.
	\item Cierre el vial con el tapón.
	\item Agite el vial para mezclar la muestra con el indicador.
	\item Observe la tonalidad que adquiere la muestra y compare con el indicador de cloro.
	\item Registre el valor de cloro obtenido con la prueba
	\item El rango de valor de concentración ideal se encuentra entre \qtyrange{1}{1.5}{\milli\gram\per\liter}
\end{enumerate}

\subsection{Responsables de la actividad}
\begin{itemize}
	\item \textbf{Ejecutado} por personal de aseguramiento de calidad.
\end{itemize}

\subsection{Acciones preventivas}
\begin{itemize}
	\item Se llevará a cabo una inspección diaria. Registrando el indicador y medida correspondiente
	\item Si la desviación se repite frecuentemente se dará curso de capacitación al personal de almacén y limpieza para que realicen eficientemente su trabajo.
	\item Si después de haber capacitado al personal de almacén y limpieza se siguen presentando desviaciones por causas injustificadas, se tomaran acciones más enérgicas con el personal por incumplimiento con sus deberes.
\end{itemize}

\subsection{Acciones correctivas}

\begin{itemize}
	\item En caso de no detectarse cloro en el agua de red municipal, el almacén cuenta con un dosificador de pulsos el cual es activado para ingresar el cloro a la línea de suministro del almacén; cabe mencionar que el agua no es utilizada para contacto con producto, solo para servicios y lavado de manos.
\end{itemize}

\subsection{Frecuencia}

\begin{itemize}
	\item Todos los días de lunes a sábado iniciando operaciones.
\end{itemize}

\begin{changelog}[simple, sectioncmd=\subsection*,label=changelog-3.20]
	\begin{version}[v=2.0, date=2023--01, author=Pablo E. Alanis]
		\item Cambio de formato;
		\item Cambios en la serialización de versiones;
		\item Correcciones ortográficas y de estilo.
	\end{version}

	\begin{version}[v=1.5, date=2022--05, author=Alonso M.]
		\item cambio de fecha;
	\end{version}	

	\shortversion{v=1.4, date=2021--05, changes=No hubo cambios}
\end{changelog}
\thispagestyle{formato-PI}
\renewcommand{\MenorVer}{0}
\renewcommand{\MayorVer}{2}
%\renewcommand{\Edit}{\MayorVer.\MenorVer}
\renewcommand{\Codigo}{HYS-22-IT}
\renewcommand{\FechaPub}{2023--01}
\renewcommand{\Titulo}{Limpieza en caso de derrame de químicos}

\section{\Titulo}\index{Limpieza!derrame de químicos, en caso de}
% \section{Procedimiento en caso de derrame de químicos}

\subsection{Objetivo}
Establecer un procedimiento de limpieza que asegure un mantenimiento y retiro apropiada de los productos químicos que han sufrido un derrame, y se realice una limpieza correcta para garantizar la eliminación de trazas de este, con el fin de evitar la contaminación cruzada.

\subsection{Alcance}
Áreas afectadas por derrame de producto químico.

\subsection{Terminología y definiciones}
\begin{description}
\defglo{limpieza}
\defglo{sanitizacion}
\end{description}

\subsection{Documentos y/o normas relacionadas}
\begin{itemize}
	\item Manual de Buenas Prácticas de Manufactura
	\item Manual de Higiene.
\end{itemize}

\subsection{Procedimiento}
\subsubsection{Instrucciones}
\begin{enumerate}
	\item Verificar la no afectación de ningún producto.
	\item Contener el derrame con el cilindro absorbente.
	\item Retirar el producto derramado con ayuda de las toallas absorbentes.
	\item Depositar en bolsa plástica
	\item Enjuagar el piso (Si fuera necesario)
	\item Eliminar la humedad con el jalador asignado
	\item Pedir al supervisor que realice una revisión visual
	\item Registrar evento en formato de derrame de químicos.
\end{enumerate}

\subsection{Responsables de la actividad}
\begin{itemize}
	\item \textbf{Ejecutado} por personal de limpieza.
\end{itemize}

\subsection{Acciones correctivas}
\begin{itemize}
	\item Cuando se presente una no conformidad en la realización del procedimiento, detectada por el supervisor o encargado del área deberá de repetir este mismo hasta que sea corregida la desviación.
	\item Se procederá según las indicaciones del departamento de calidad del cliente.
\end{itemize}

\subsection{Frecuencia}
\begin{itemize}
	\item Cada que ocurra un evento de derrame de químicos
\end{itemize}

\begin{changelog}[title=Registro de cambios,simple, sectioncmd=\subsection*,label=changelog-\thesection-\MayorVer.\MenorVer2]
	\begin{version}[v=\MayorVer.\MenorVer, date=2023--01, author=Pablo E. Alanis]
		\item Cambio de formato;
		\item Cambios en la serialización de versiones;
		\item Correcciones ortográficas y de estilo.
	\end{version}

	\begin{version}[v=1.5, date=2022--05, author=Alonso M.]
		\item cambio de fecha;
	\end{version}	

	\shortversion{v=1.4, date=2021--05, changes=No hubo cambios}
\end{changelog}
\thispagestyle{formato-PI}
\renewcommand{\MenorVer}{0}
\renewcommand{\MayorVer}{2}
%\renewcommand{\Edit}{\MayorVer.\MenorVer}
\renewcommand{\Codigo}{HYS-23-IT}
\renewcommand{\FechaPub}{2023--01}
\renewcommand{\Titulo}{Limpieza y desinfección de cisterna}

\section{\Titulo}\index{Limpieza!cisterna, de}
% \section{Limpieza y desinfección de cisterna}

\subsection{Objetivo}
Establecer un método para la correcta limpieza y desinfección de cisterna.

\subsection{Alcance}
Aplica a cisternas de agua.

\subsection{Terminología y definiciones}
\begin{description}
\defglo{limpieza}
\defglo{sanitizacion}
\defglo{cisterna}
\end{description}

\subsection{Procedimiento}
\begin{itemize}
	\item[\textbf{Aviso:}] \gls{RDF} NO CUENTA CON CISTERNA DE AGUA dentro de sus operaciones.
	\item[\textbf{Aviso:}] \gls{RDF} NO realiza actividades de LAVADO Y DESINFECCION DE CISTERNA DE AGUA como práctica común dentro de sus procedimientos durante la operación.
	\item[\textbf{Aviso:}] \gls{RDF} NO cuenta con NINGUN tipo de almacenamiento de Agua dentro de las instalaciones del almacén.
	\item[\textbf{Aviso:}] \gls{RDF} utiliza agua suministrada por la red municipal de Monterrey, esta agua no tiene contacto directo en ningún momento con el producto.
\end{itemize}

\thispagestyle{formato-PI}
\renewcommand{\MenorVer}{0}
\renewcommand{\MayorVer}{2}
%\renewcommand{\Edit}{\MayorVer.\MenorVer}
\renewcommand{\Codigo}{HYS-24-IT}
\renewcommand{\FechaPub}{2023--01}
\renewcommand{\Titulo}{Cloración de agua potable}

\section{\Titulo}\index{Cloración!agua potable, de}
% \section{Cloración de agua potable}

\subsection{Objetivo}
Establecer las instrucciones de trabajo para la cloración del agua potable almacenada en el tinaco.

\subsection{Alcance}
\begin{itemize}
	\item Personal de mantenimiento;
	\item personal de aseguramiento de calidad.
\end{itemize}

\subsection{Terminología y definiciones}
\begin{description}
\defglo{cloración}
\end{description}

\subsection{Documentos y/o normas relacionadas}
\begin{itemize}
	\item Codex Alimentarius
	\item \BitClor
\end{itemize}

\subsection{Procedimiento}
\subsubsection{Precauciones de seguridad}
\begin{itemize}
	\item El indicador empleado para la determinación de cloro, 2-metilanilina, es cancerígeno. Se deben de seguir las indicaciones del MSDS y evitar la exposición directa a este compuesto.
\end{itemize}

\subsubsection{Materiales}
\begin{itemize}
	\item \BitClor.
	% \item Bitácora de verificación de cloración de agua potable. %TODO: ponerlo algun dia
\end{itemize}

\subsubsection{Instrucciones}

\begin{enumerate}
	\item Limpiar tapa del tinaco;
	\item Abrir tapa del tinaco y colocarla sobre una superficie limpia;
	\item Colocar las pastillas de cloro requeridas según las especificaciones del fabricante.
	\item Cerrar tinaco y limpiar nuevamente la tapa;
	\item Anotar en bitácora de cloración de agua.
\end{enumerate}

\subsection{Responsables de la actividad}

\begin{itemize}
	\item \textbf{Ejecutado} por personal de mantenimiento;
	\item \textbf{Monitoreado} por personal de calidad;
	\item \textbf{Verificado} por personal de calidad.
\end{itemize}

\subsection{Acciones preventivas}

\begin{itemize}
	\item Se llevará a cabo una determinación cualitativa del cloro libre presente en el suministro de agua empleando 2-metilanilina (2-amino-1-metilbenceno) dos veces al mes y se registrará la concentración de cloro en la \emph{bitácora de verificación de cloración de agua potable.}
\end{itemize}

\subsection{Acciones correctivas}

\begin{itemize}
	\item Si se detecta que la concentración de cloro es inferior al punto de vire de la 2-metilanilina, debe de reportarse en la \emph{bitácora de verificación de cloración de agua potable} y reportarse a la gerencia.
\end{itemize}

\subsection{Frecuencia}

\begin{itemize}
	\item Mensualmente.
	% \item[Verificación] mensualmente.
\end{itemize}

\subsection{Historial de modificaciones}

\begin{changelog}[simple, sectioncmd=\subsection*,label=changelog-3.22]
	\shortversion{v=1.0, author=Pablo E. Alanis, date=2022--12, changes=Primera edición}
\end{changelog}
\renewcommand{\MenorVer}{0}
\renewcommand{\MayorVer}{2}
%\renewcommand{\Edit}{\MayorVer.\MenorVer}
\renewcommand{\Codigo}{HYS-25-L}
\renewcommand{\FechaPub}{2023--01}
\renewcommand{\Titulo}{Lista maestra de productos químicos}

\section{\Titulo}
% \section{Lista maestra de productos químicos}

\subsection{Sanidad}

\subsubsection{Liquid-K}


\subsubsection{\textit{Kool Klene}}


\subsubsection{Hands free}


\subsubsection{Gelshand}


\subsubsection{Whisper V}


\subsection{Montacargas}

\begin{itemize}
	\item Agua destilada
	\item Aceite hidráulico
	\item Grasa grado alimenticio
\textbf{Nota:} En nuestro caso somos un almacén y nuestros lavados son en superficies de no contacto con producto y el producto no se encuentra expuesto (cuenta con empaque primario y secundario) por lo cual no tenemos programa de rotación de productos químicos de limpieza. Contamos con un programa de limpieza donde los periodos de lavado profundo no son cortos por lo cual no requerimos la rotación de productos químicos.
\end{itemize}

\renewcommand{\MayorVer}{1}
\renewcommand{\MenorVer}{0}
\renewcommand{\Codigo}{HyS-G}
\renewcommand{\FechaPub}{2023--01}
%\renewcommand{\Edit}{01}
\renewcommand{\Titulo}{Programa de Higiene y Sanidad --- Glosario}

\printacronyms
\printglossary 
\renewcommand{\MayorVer}{1}
\renewcommand{\MenorVer}{0}
\renewcommand{\Codigo}{HyS-I}
\renewcommand{\FechaPub}{2023--01}
%\renewcommand{\Edit}{01}
\renewcommand{\Titulo}{Indice alfabético}
\printindex
\end{document}